% Luc Destombe --- licence CC BY-NC-SA 4.0
% https://creativecommons.org/licenses/by-nc-sa/4.0/deed.fr
% Fichier autonome : compiler avec pdflatex, deux fois (signets, nombre de pages).
\documentclass[11pt,a4paper]{article}
\makeatletter
\newif\iflycee@unecolonne
\lycee@unecolonnetrue
\newif\iflycee@palatino
\lycee@palatinotrue

\RequirePackage[utf8]{inputenc}
\RequirePackage[T1]{fontenc}
\RequirePackage[french]{babel}
\RequirePackage{etoolbox}

\newcommand{\ShorthandsOff}{\@ifundefined{shorthandoff}{}{\shorthandoff{;:!?}}}
\newcommand{\ShorthandsOn}{\@ifundefined{shorthandon}{}{\shorthandon{;:!?}}}
\ShorthandsOff
\AtBeginDocument{\ShorthandsOn}
\AtBeginEnvironment{tikzpicture}{\ShorthandsOff}

\RequirePackage{geometry}
\RequirePackage{fancyhdr}
\RequirePackage{lastpage}
\RequirePackage{multicol}
\RequirePackage{array}
\RequirePackage{enumitem}
\RequirePackage{xcolor}
\RequirePackage{needspace}

\RequirePackage{amsmath,amssymb}
\RequirePackage{mathpazo}
\DeclareMathSizes{10.95}{10.95}{8}{5.5}
\begingroup\catcode`\_=\active \catcode`\^=\active
\gdef_#1{\sb{\scriptscriptstyle #1}}
\gdef^#1{\sp{\scriptscriptstyle\lycee@moinsepais #1}}
\endgroup
\newcommand\lycee@moins{\mathbin{%
  \setbox\z@\hbox{$\m@th\scriptscriptstyle\lycee@moinsorig$}%
  \dimen@=.55\wd\z@ \advance\dimen@-.5pt
  \hbox to.75\wd\z@{\hss$\m@th\scriptscriptstyle\vcenter{\hbox{%
    \tikz\draw[line width=.5pt,line cap=round](0,0)--(\the\dimen@,0);}}$\hss}}}
\begingroup\lccode`\~=`\- \lowercase{\endgroup\let~}\lycee@moins
\newcommand\lycee@moinsepais{\mathcode`\-="8000 }
\AtBeginDocument{\mathchardef\lycee@moinsorig=\mathcode`\-\relax}
\AtBeginDocument{\catcode`\_=12 \mathcode`\_="8000
  \catcode`\^=12 \mathcode`\^="8000 }
\newcommand\lycee@exposants{%
  \fontdimen13\textfont2=.48\fontdimen6\textfont2
  \fontdimen14\textfont2=.48\fontdimen6\textfont2
  \fontdimen15\textfont2=.48\fontdimen6\textfont2
  \fontdimen13\scriptfont2=.48\fontdimen6\scriptfont2
  \fontdimen14\scriptfont2=.48\fontdimen6\scriptfont2
  \fontdimen15\scriptfont2=.48\fontdimen6\scriptfont2}
\every@math@size\expandafter{\the\every@math@size\lycee@exposants}
\newlength{\retraitcalculs}
\setlength{\retraitcalculs}{2em}

\RequirePackage{tikz}
\usetikzlibrary{arrows.meta,calc,babel,shapes.misc}
\RequirePackage[most]{tcolorbox}
\tcbuselibrary{skins,breakable}

\definecolor{coulDef}{HTML}{1F4E79}
\definecolor{coulProp}{HTML}{7030A0}
\definecolor{coulRem}{HTML}{C55A11}
\definecolor{coulEx}{HTML}{2E7D32}
\definecolor{coulExo}{HTML}{B03A2E}
\definecolor{coulPart}{HTML}{1F4E79}
\definecolor{coulMeth}{HTML}{1B6E4F}
\definecolor{coulCourbe}{HTML}{0B5ED7}
\definecolor{coulCourbeB}{HTML}{C00000}
\definecolor{coulCourbeC}{HTML}{2E7D32}
\definecolor{coulGrille}{HTML}{8C8C8C}
\definecolor{coulRep}{HTML}{1B6E4F}
\definecolor{coulBar}{HTML}{2E7D32}

\newcommand{\R}{\mathbb{R}}
\newcommand{\Rs}{\mathbb{R}^{*}}
\renewcommand{\le}{\leqslant}
\renewcommand{\ge}{\geqslant}
\renewcommand{\approx}{\simeq}
\newcommand{\vect}[1]{\overrightarrow{#1}}
\newcommand{\norme}[1]{\left\lVert #1 \right\rVert}
\newcommand{\intervalle}[2]{\left[#1\,;#2\right]}
\RequirePackage{cancel}
\renewcommand{\CancelColor}{\color{coulExo}}
\newcommand{\fois}[1]{\times{\color{coulExo}#1}}
\newcommand{\barre}[1]{\cancel{#1}}

\tikzset{grille/.style={coulGrille,line width=0.4pt}}
\newcommand{\quadrillage}[4]{\draw[grille,step=1] (#1,#2) grid (#3,#4);}
\tikzset{
  croix/.style={cross out,draw,line width=0.6pt,inner sep=0pt,minimum size=4pt},
  croixdroite/.style={croix,rotate=45,line width=0.8pt},
}
\newcommand{\pt}[3]{\path (#1) node[croix]{} node[#2,font=\small,inner sep=2.5pt]{$#3$};}

\newcommand{\lycee@repere}[4]{%
  \draw[grille,step=1] (#1,#2) grid (#3,#4);
  \draw[-{Stealth[length=2mm]},thick] (#1,0)--({#3+0.4},0) node[below left=-1pt]{$x$};
  \draw[-{Stealth[length=2mm]},thick] (0,#2)--(0,{#4+0.4}) node[below left=-1pt]{$y$};
  \node[below left,font=\scriptsize,inner sep=1.5pt] at (0,0) {$0$};}
\newcommand{\lycee@gradx}[1]{\foreach \k/\l in {#1}{%
  \draw (\k,0.07)--(\k,-0.07) node[below,font=\scriptsize,inner sep=1.5pt]{$\l$};}}
\newcommand{\lycee@grady}[1]{\foreach \k/\l in {#1}{%
  \draw (0.07,\k)--(-0.07,\k) node[left,font=\scriptsize,inner sep=1.5pt]{$\l$};}}
\newcommand{\lycee@intff}[2]{\left[#1\,;#2\right]}
\newcommand{\lycee@intfo}[2]{\left[#1\,;#2\right[}
\newcommand{\lycee@intof}[2]{\left]#1\,;#2\right]}
\newcommand{\lycee@intoo}[2]{\left]#1\,;#2\right[}
\newcommand{\lycee@ens}[1]{\left\{#1\right\}}
\AtBeginDocument{%
  \providecommand{\repere}{\lycee@repere}%
  \providecommand{\gradx}{\lycee@gradx}%
  \providecommand{\grady}{\lycee@grady}%
  \providecommand{\Cf}{\mathcal{C}\sb{\scriptscriptstyle f}}%
  \providecommand{\Cg}{\mathcal{C}\sb{\scriptscriptstyle g}}%
  \providecommand{\intff}{\lycee@intff}%
  \providecommand{\intfo}{\lycee@intfo}%
  \providecommand{\intof}{\lycee@intof}%
  \providecommand{\intoo}{\lycee@intoo}%
  \providecommand{\ens}{\lycee@ens}}

\newlist{questions}{enumerate}{3}
\setlist[questions,1]{label=\arabic*\textdegree),leftmargin=*,itemsep=2pt,topsep=3pt}
\setlist[questions,2]{label=\alph*),leftmargin=*,itemsep=1pt,topsep=2pt}
\setlist[questions,3]{label=\roman*),leftmargin=*,itemsep=1pt,topsep=1pt}
\setlist[itemize]{label=\textbullet,leftmargin=*,itemsep=2pt,topsep=3pt}

\newcommand{\besoinplace}[1]{\par\penalty\@M\begingroup
  \dimen@=#1\relax\dimen@ii=\pagegoal\advance\dimen@ii-\pagetotal
  \ifdim\dimen@>\dimen@ii\ifdim\dimen@ii>\z@\vfil\fi\break\fi\endgroup}

\newcommand{\lycee@pied}{\thepage/\pageref{LastPage}}
\newcommand{\entete}[2]{%
  \pagestyle{fancy}\fancyhf{}%
  \fancyhead[L]{\footnotesize\color{coulPart}\textbf{#1}}%
  \fancyhead[R]{\footnotesize\color{coulPart}\textbf{#2}}%
  \fancyfoot[C]{\footnotesize\lycee@pied}%
  \renewcommand{\headrulewidth}{0.4pt}%
  \renewcommand{\headrule}{\hbox to\headwidth{%
    \color{coulPart!40}\leaders\hrule height \headrulewidth\hfill}}}

\RequirePackage{ccicons}
\newcommand{\lycee@licence}{{\scriptsize\color{gray}%
  \copyright\ Luc Destombe --- \ccbyncsa\ CC BY-NC-SA 4.0}}
\newif\iflycee@licence \lycee@licencetrue
\newcommand{\sanslicence}{\lycee@licencefalse}
\AtBeginDocument{\iflycee@licence\AddToHookNext{shipout/foreground}{%
  \put(0,0){\rlap{\kern\dimexpr1in+\hoffset+\oddsidemargin+\textwidth\relax
    \llap{\raisebox{-\dimexpr1in+\voffset+\topmargin+\headheight+\headsep
      +\textheight+\footskip\relax}{\lycee@licence}}}}}\fi}

\newcommand{\formulesserrees}{%
  \setlength{\abovedisplayskip}{3pt}\setlength{\belowdisplayskip}{3pt}%
  \setlength{\abovedisplayshortskip}{2pt}\setlength{\belowdisplayshortskip}{3pt}}

\setlength{\parindent}{0pt}

\RequirePackage[implicit=false,hidelinks,pdfencoding=auto,psdextra,bookmarksopen,
  bookmarksopenlevel=1,pdfcreator={},pdfproducer={}]{hyperref}
\RequirePackage{bookmark}
\hypersetup{pdfauthor={Luc Destombe},pdfkeywords={CC BY-NC-SA 4.0}}
\newcounter{lycee@signet}
\newcommand{\signet}[3][]{\stepcounter{lycee@signet}%
  \edef\lycee@dest{\if\relax\detokenize{#1}\relax
    signet.\arabic{lycee@signet}\else#1\fi}%
  \hypertarget{\lycee@dest}{}\bookmark[level=#2,dest=\lycee@dest]{#3}}
\pdfstringdefDisableCommands{%
  \def\R{R}\def\vect#1{#1}\def\Cf{Cf}\def\Cg{Cg}\def\fois#1{×#1}%
  \def\barre#1{#1}\def\intervalle#1#2{[#1;#2]}\def\intff#1#2{[#1;#2]}%
  \def\textsuperscript#1{#1}\def\dfrac#1#2{#1/#2}\def\frac#1#2{#1/#2}%
  \def\vec#1{#1⃗}}%
\RequirePackage[official]{eurosym}

\tcbset{exercice corrige/.style={
  enhanced,breakable,before upper=\raggedright,
  colback=white,colframe=coulExo,
  boxrule=0.8pt,arc=2pt,
  left=6pt,right=6pt,top=8pt,bottom=5pt,
  before skip=9pt,after skip=4pt,
  coltitle=white,fonttitle=\bfseries\small,
  attach boxed title to top left={xshift=8pt,yshift=-\tcboxedtitleheight/2},
  boxed title style={colback=coulExo,arc=2pt,boxrule=0pt}}}
\newcounter{exo}
\newtcolorbox{exercice}[1][]{exercice corrige,
  phantom={\signet[exercice-\arabic{exo}]{2}{Exercice \arabic{exo}}},
  title={Exercice \theexo\ifx\relax#1\relax\else{} --- #1\fi}}
\newenvironment{exo}[1][\relax]{\refstepcounter{exo}\begin{exercice}[#1]}{\end{exercice}}

\geometry{top=1.8cm,bottom=1cm,left=1cm,right=1cm,headheight=14pt,footskip=0.6cm}
\renewcommand{\lycee@pied}{\thepage}

\renewcommand{\theexo}{\ifnum\value{exo}<10 0\fi\arabic{exo}}
\tcbset{exercice corrige/.style={
  enhanced,breakable,before upper=\raggedright,
  colback=white,colframe=coulExo,
  boxrule=0.9pt,arc=2pt,
  left=8pt,right=8pt,top=8pt,bottom=6pt,
  before skip=8pt,after skip=6pt,
  coltitle=white,fonttitle=\bfseries,
  attach boxed title to top left={xshift=8pt,yshift=-\tcboxedtitleheight/2},
  boxed title style={colback=coulExo,arc=2pt,boxrule=0pt}}}

\newcommand{\entetefiche}[2]{%
  \begin{center}
    \begin{tcolorbox}[enhanced,width=0.92\linewidth,colback=white,colframe=coulPart,
        boxrule=1.4pt,arc=3pt,halign=center,top=6pt,bottom=6pt]
      {\large\bfseries\color{coulPart} CORRIGÉ --- FICHE D'EXERCICES}\\[3pt]
      {\normalsize\bfseries #1}\\[2pt]
      {\normalsize #2}
    \end{tcolorbox}
  \end{center}}

\newcounter{partie}
\newcommand{\partie}[1]{%
  \stepcounter{partie}%
  \besoinplace{8\baselineskip}\vspace{4pt}%
  \begin{tcolorbox}[enhanced,colback=coulPart!8,colframe=coulPart,boxrule=0pt,
      leftrule=3pt,arc=0pt,left=6pt,right=4pt,top=3pt,bottom=3pt,
      before skip=6pt,after skip=2pt]
    \bfseries\color{coulPart}\Roman{partie} --- #1\signet{1}{\Roman{partie} --- #1}
  \end{tcolorbox}}
\newcommand{\partiesansnum}[1]{%
  \besoinplace{8\baselineskip}\vspace{4pt}%
  \begin{tcolorbox}[enhanced,colback=coulPart!8,colframe=coulPart,boxrule=0pt,
      leftrule=3pt,arc=0pt,left=6pt,right=4pt,top=3pt,bottom=3pt,
      before skip=6pt,after skip=2pt]
    \bfseries\color{coulPart}#1\signet{1}{#1}
  \end{tcolorbox}}

\setlist[questions,1]{label=\arabic*\textdegree),leftmargin=*,itemsep=3pt,topsep=3pt}
\setlist[questions,2]{label=\alph*),leftmargin=*,itemsep=2pt,topsep=2pt}
\setlist[questions,3]{label=\roman*),leftmargin=*,itemsep=1pt,topsep=1pt}
\newlist{expr}{enumerate}{1}
\setlist[expr]{label={\Alph*\ $=$},leftmargin=*,itemsep=2pt,topsep=3pt}

\newcommand{\res}[1]{{\color{coulRep}#1}}
\newcommand{\rem}[1]{\par\smallskip{\small\itshape\color{black!60}$\rightarrow$ #1}}
\newenvironment{calculs}{\par\smallskip\noindent$\displaystyle\begin{aligned}[t]}%
  {\end{aligned}$\par\smallskip}
\newcommand{\justif}[1]{\quad\text{\small\itshape\color{black!60}#1}}
\newcommand{\cas}[1]{\par\smallskip\noindent\textbf{\color{coulExo}#1}\nobreak}
\newcommand{\explique}[2]{\par\smallskip\noindent
  \begin{minipage}[t]{0.6\linewidth}\raggedright #1\end{minipage}\hfill
  \begin{minipage}[t]{0.37\linewidth}\raggedright\leavevmode
    {\small\itshape\color{black!60}#2}\end{minipage}\par}

\setlength{\parskip}{2pt}

\makeatother
\entete{Chapitre 4 --- Dérivation --- corrigé}{Première spécialité}
\usepackage{tkz-tab}
\let\casligne\cas
\renewcommand{\cas}[1]{\casligne{#1}\par\nopagebreak}
\setlength{\parindent}{0pt}

\begin{document}

\entetefiche{Chapitre 4 --- Dérivation}{Première spécialité mathématiques}

\partie{Limite en zéro d'une fonction}

\begin{exo}[Conjecturer avec un tableau de valeurs]
\cas{1.}
\explique{On remplace $x$ par chaque valeur dans $f(x)=\dfrac{x^{2}-3x}{x}$ ; on détaille
la première.}{$0$ est une valeur interdite : on s'en approche sans l'atteindre}
\begin{calculs}
f(-0{,}1) &= \dfrac{(-0{,}1)^{2}-3\times(-0{,}1)}{-0{,}1} &&\justif{on remplace $x$ par $-0{,}1$}\\
f(-0{,}1) &= \dfrac{0{,}01+0{,}3}{-0{,}1} &&\justif{un carré est positif ; $-3\times(-0{,}1)=+0{,}3$}\\
f(-0{,}1) &= \res{-3{,}1} &&\justif{$0{,}31$ divisé par $-0{,}1$}
\end{calculs}
\explique{De même pour les trois autres valeurs :}{calculatrice possible}
\begin{center}
\renewcommand{\arraystretch}{1.3}
\begin{tabular}{|c|c|c|c|c|}
\hline
$x$ & $-0{,}1$ & $-0{,}01$ & $0{,}01$ & $0{,}1$\\ \hline
$f(x)$ & $\res{-3{,}1}$ & $\res{-3{,}01}$ & $\res{-2{,}99}$ & $\res{-2{,}9}$\\ \hline
\end{tabular}
\end{center}
\cas{2.}
\explique{Quand $x$ s'approche de $0$, d'un côté comme de l'autre, $f(x)$ s'approche de
$-3$ : la limite semble être $\res{-3}$.}{$-3{,}1$ puis $-3{,}01$ ; $-2{,}9$ puis
$-2{,}99$ : l'écart à $-3$ est divisé par $10$}
\cas{3.}
\begin{calculs}
f(x) &= \dfrac{x^{2}-3x}{x} &&\justif{on recopie l'expression}\\
f(x) &= \dfrac{\barre{x}(x-3)}{\barre{x}} &&\justif{on factorise le numérateur par $x$}\\
f(x) &= x-3 &&\justif{on simplifie : $x\ne 0$}\\
\lim_{x\to 0} f(x) &= \res{-3} &&\justif{on remplace $x$ par $0$ dans $x-3$}
\end{calculs}
\explique{On retrouve la valeur conjecturée.}{le tableau donne une idée, le calcul la
prouve}
\rem{On ne remplace pas $x$ par $0$ dans l'expression de départ : on obtiendrait
$\frac{0}{0}$, qui n'a pas de sens. On simplifie d'abord par $x$.}
\end{exo}

\begin{exo}[Calculer des limites en zéro]
\explique{Dans chaque numérateur, $x$ est en facteur : on factorise par $x$, on simplifie,
puis on remplace $x$ par $0$.}{méthode 1}
\cas{a)}
\begin{calculs}
f(x) &= \dfrac{x^{2}+4x}{x} &&\justif{on recopie l'expression}\\
f(x) &= \dfrac{\barre{x}(x+4)}{\barre{x}} &&\justif{on factorise le numérateur par $x$}\\
f(x) &= x+4 &&\justif{on simplifie : $x\ne 0$}\\
\lim_{x\to 0} f(x) &= \res{4} &&\justif{on remplace $x$ par $0$ dans $x+4$}
\end{calculs}
\cas{b)}
\begin{calculs}
g(x) &= \dfrac{2x^{2}-6x}{x} &&\\
g(x) &= \dfrac{\barre{x}(2x-6)}{\barre{x}} &&\justif{on factorise par $x$}\\
g(x) &= 2x-6 &&\\
\lim_{x\to 0} g(x) &= \res{-6} &&\justif{$2x$ tend vers $0$}
\end{calculs}
\cas{c)}
\begin{calculs}
h(x) &= \dfrac{4x^{3}+x}{x} &&\\
h(x) &= \dfrac{\barre{x}(4x^{2}+1)}{\barre{x}} &&\justif{$x=x\times 1$ : le $1$ ne s'oublie pas}\\
h(x) &= 4x^{2}+1 &&\\
\lim_{x\to 0} h(x) &= \res{1} &&\justif{$4x^{2}$ tend vers $0$}
\end{calculs}
\cas{d)}
\begin{calculs}
k(x) &= \dfrac{3x^{3}-x^{2}+8x}{x} &&\\
k(x) &= \dfrac{\barre{x}(3x^{2}-x+8)}{\barre{x}} &&\justif{chaque terme contient $x$}\\
k(x) &= 3x^{2}-x+8 &&\\
\lim_{x\to 0} k(x) &= \res{8} &&\justif{$3x^{2}$ et $-x$ tendent vers $0$}
\end{calculs}
\rem{Erreur fréquente en c) : écrire $4x^{3}+x=x\times 4x^{2}$ en oubliant le $1$. On
vérifie une factorisation en redéveloppant : $x(4x^{2}+1)=4x^{3}+x$.}
\end{exo}

\begin{exo}[Avec une identité remarquable]
\explique{On ne peut pas factoriser tout de suite : on développe d'abord le carré, la
constante disparaît, puis on factorise par $x$.}{méthode 1, exemple c)}
\cas{a)}
\begin{calculs}
f(x) &= \dfrac{(1+x)^{2}-1}{x} &&\justif{on recopie l'expression}\\
f(x) &= \dfrac{1+2x+x^{2}-1}{x} &&\justif{$(a+b)^{2}=a^{2}+2ab+b^{2}$}\\
f(x) &= \dfrac{\barre{x}(2+x)}{\barre{x}} &&\justif{$1$ et $-1$ s'annulent}\\
f(x) &= 2+x &&\\
\lim_{x\to 0} f(x) &= \res{2} &&\justif{on remplace $x$ par $0$}
\end{calculs}
\cas{b)}
\begin{calculs}
g(x) &= \dfrac{(5+x)^{2}-25}{x} &&\\
g(x) &= \dfrac{25+10x+x^{2}-25}{x} &&\justif{double produit : $2\times 5\times x=10x$}\\
g(x) &= \dfrac{\barre{x}(10+x)}{\barre{x}} &&\justif{$25$ et $-25$ s'annulent}\\
g(x) &= 10+x &&\\
\lim_{x\to 0} g(x) &= \res{10} &&
\end{calculs}
\cas{c)}
\begin{calculs}
h(x) &= \dfrac{(x-2)^{2}-4}{x} &&\\
h(x) &= \dfrac{x^{2}-4x+4-4}{x} &&\justif{$(a-b)^{2}=a^{2}-2ab+b^{2}$}\\
h(x) &= \dfrac{\barre{x}(x-4)}{\barre{x}} &&\justif{$4$ et $-4$ s'annulent}\\
h(x) &= x-4 &&\\
\lim_{x\to 0} h(x) &= \res{-4} &&
\end{calculs}
\cas{d)}
\begin{calculs}
k(x) &= \dfrac{(3-x)^{2}-9}{x} &&\\
k(x) &= \dfrac{9-6x+x^{2}-9}{x} &&\justif{$(a-b)^{2}$, avec $a=3$ et $b=x$}\\
k(x) &= \dfrac{\barre{x}(-6+x)}{\barre{x}} &&\justif{$9$ et $-9$ s'annulent}\\
k(x) &= x-6 &&\\
\lim_{x\to 0} k(x) &= \res{-6} &&
\end{calculs}
\rem{Erreur fréquente : écrire $(1+x)^{2}=1+x^{2}$. Le double produit est justement le
terme qui donne la limite.}
\end{exo}

\begin{exo}[Limites quand $h$ tend vers $0$]
\explique{Même méthode avec la lettre $h$ à la place de $x$ : c'est sous cette forme que
ces limites reviennent dans le calcul d'un nombre dérivé (partie II).}{$h$ s'approche de
$0$ sans jamais valoir $0$}
\cas{a)}
\begin{calculs}
\dfrac{h^{2}+6h}{h} &= \dfrac{\barre{h}(h+6)}{\barre{h}} &&\justif{on factorise par $h$}\\
\dfrac{h^{2}+6h}{h} &= h+6 &&\justif{on simplifie : $h\ne 0$}\\
\lim_{h\to 0}\dfrac{h^{2}+6h}{h} &= \res{6} &&\justif{on remplace $h$ par $0$}
\end{calculs}
\cas{b)}
\begin{calculs}
\dfrac{(2+h)^{2}-4}{h} &= \dfrac{4+4h+h^{2}-4}{h} &&\justif{identité remarquable}\\
\dfrac{(2+h)^{2}-4}{h} &= \dfrac{\barre{h}(4+h)}{\barre{h}} &&\justif{$4$ et $-4$ s'annulent}\\
\dfrac{(2+h)^{2}-4}{h} &= 4+h &&\\
\lim_{h\to 0}\dfrac{(2+h)^{2}-4}{h} &= \res{4} &&
\end{calculs}
\cas{c)}
\begin{calculs}
\dfrac{3(1+h)^{2}-3}{h} &= \dfrac{3(1+2h+h^{2})-3}{h} &&\justif{on développe d'abord le carré}\\
\dfrac{3(1+h)^{2}-3}{h} &= \dfrac{3+6h+3h^{2}-3}{h} &&\justif{puis on multiplie par $3$}\\
\dfrac{3(1+h)^{2}-3}{h} &= \dfrac{\barre{h}(6+3h)}{\barre{h}} &&\\
\dfrac{3(1+h)^{2}-3}{h} &= 6+3h &&\\
\lim_{h\to 0}\dfrac{3(1+h)^{2}-3}{h} &= \res{6} &&\justif{$3h$ tend vers $0$}
\end{calculs}
\rem{En c), le $3$ multiplie tout le carré : on calcule le carré d'abord (priorité des
puissances), puis on multiplie par $3$.}
\end{exo}

\partie{Nombre dérivé}

\begin{exo}[Nombre dérivé par la définition]
\explique{Pour chaque fonction : on calcule l'image de $a$, puis celle de $a+h$, le taux
$T(h)=\dfrac{f(a+h)-f(a)}{h}$, et enfin sa limite quand $h$ tend vers $0$.}{méthode 2}
\cas{a)}
\begin{calculs}
f(3) &= 9 &&\justif{on calcule d'abord $f(3)=3^{2}$}\\
f(3+h) &= (3+h)^{2} &&\justif{on remplace $x$ par $3+h$}\\
f(3+h) &= 9+6h+h^{2} &&\justif{$(a+b)^{2}=a^{2}+2ab+b^{2}$}\\
T(h) &= \dfrac{9+6h+h^{2}-9}{h} &&\justif{$T(h)=\frac{f(3+h)-f(3)}{h}$}\\
T(h) &= \dfrac{\barre{h}(6+h)}{\barre{h}} &&\justif{$9$ et $-9$ s'annulent}\\
T(h) &= 6+h &&\justif{on simplifie : $h\ne 0$}\\
f'(3) &= \lim_{h\to 0}\,(6+h) &&\justif{$f'(3)$ est la limite de $T(h)$}\\
f'(3) &= \res{6} &&\justif{on remplace $h$ par $0$}
\end{calculs}
\cas{b)}
\begin{calculs}
g(1) &= 1+2 &&\justif{on calcule d'abord $g(1)$}\\
g(1) &= 3 &&\\
g(1+h) &= (1+h)^{2}+2(1+h) &&\justif{on remplace $x$ par $1+h$, partout}\\
g(1+h) &= 1+2h+h^{2}+2+2h &&\\
g(1+h) &= h^{2}+4h+3 &&\\
T(h) &= \dfrac{h^{2}+4h+3-3}{h} &&\justif{on soustrait $g(1)=3$}\\
T(h) &= \dfrac{\barre{h}(h+4)}{\barre{h}} &&\\
T(h) &= h+4 &&\\
g'(1) &= \lim_{h\to 0}\,(h+4) &&\\
g'(1) &= \res{4} &&
\end{calculs}
\cas{c)}
\begin{calculs}
k(1) &= -1+5 &&\justif{$-1^{2}=-1$ : le carré porte sur $1$}\\
k(1) &= 4 &&\\
k(1+h) &= -(1+h)^{2}+5(1+h) &&\justif{le \og$-$\fg{} porte sur tout le carré}\\
k(1+h) &= -(1+2h+h^{2})+5+5h &&\\
k(1+h) &= -1-2h-h^{2}+5+5h &&\justif{on enlève les parenthèses : tous les signes changent}\\
k(1+h) &= -h^{2}+3h+4 &&\\
T(h) &= \dfrac{-h^{2}+3h+4-4}{h} &&\\
T(h) &= \dfrac{\barre{h}(-h+3)}{\barre{h}} &&\\
T(h) &= -h+3 &&\\
k'(1) &= \lim_{h\to 0}\,(-h+3) &&\\
k'(1) &= \res{3} &&
\end{calculs}
\cas{d)}
\begin{calculs}
m(-1) &= 2\times(-1)^{2}-3 &&\justif{$(-1)^{2}=1$}\\
m(-1) &= -1 &&\\
m(-1+h) &= 2(-1+h)^{2}-3 &&\\
m(-1+h) &= 2(1-2h+h^{2})-3 &&\justif{$(-1+h)^{2}=1-2h+h^{2}$}\\
m(-1+h) &= 2h^{2}-4h-1 &&\justif{$2-4h+2h^{2}-3$}\\
T(h) &= \dfrac{2h^{2}-4h-1-(-1)}{h} &&\justif{on soustrait $m(-1)=-1$ : parenthèses}\\
T(h) &= \dfrac{\barre{h}(2h-4)}{\barre{h}} &&\justif{$-1+1=0$}\\
T(h) &= 2h-4 &&\\
m'(-1) &= \lim_{h\to 0}\,(2h-4) &&\\
m'(-1) &= \res{-4} &&
\end{calculs}
\rem{Deux pièges : le \og$-$\fg{} devant un carré, $-(1+h)^{2}=-1-2h-h^{2}$, et la
soustraction d'une image négative, $-(-1)=+1$. Contrôle : au numérateur de $T(h)$, les
constantes doivent toujours s'annuler (tous les termes restants contiennent $h$) ;
sinon, il y a une erreur de calcul.}
\end{exo}

\begin{exo}[Le taux de variation est donné]
\explique{Le nombre dérivé est la limite de $T(h)$ quand $h$ tend vers $0$. En a) et b),
$T(h)$ est déjà simplifié : il suffit de remplacer $h$ par $0$.}{définition 2}
\cas{a)}
\begin{calculs}
f'(1) &= \lim_{h\to 0}\,(3h+7) &&\justif{$f'(1)$ est la limite de $T(h)$}\\
f'(1) &= \res{7} &&\justif{$3h$ tend vers $0$}
\end{calculs}
\cas{b)}
\begin{calculs}
g'(4) &= \lim_{h\to 0}\,(h^{2}-h-2) &&\\
g'(4) &= \res{-2} &&\justif{$h^{2}$ et $-h$ tendent vers $0$}
\end{calculs}
\explique{Le nombre $4$ n'apparaît pas dans le calcul : il indique seulement en quel nombre
on dérive.}{}
\cas{c)}
\explique{Attention : on connaît ici la différence $k(-1+h)-k(-1)$, pas encore le taux. Il
faut d'abord diviser par $h$.}{$T(h)=\frac{k(-1+h)-k(-1)}{h}$}
\begin{calculs}
T(h) &= \dfrac{5h-2h^{2}}{h} &&\justif{on divise la différence par $h$}\\
T(h) &= \dfrac{\barre{h}(5-2h)}{\barre{h}} &&\justif{on factorise par $h$}\\
T(h) &= 5-2h &&\\
k'(-1) &= \lim_{h\to 0}\,(5-2h) &&\\
k'(-1) &= \res{5} &&
\end{calculs}
\rem{Erreur fréquente en c) : remplacer $h$ par $0$ dans $5h-2h^{2}$, ce qui donne $0$. La
différence tend toujours vers $0$ ; c'est le taux qui donne le nombre dérivé.}
\end{exo}

\begin{exo}[Avec des fractions]
\cas{1.}
\explique{On met les deux fractions au même dénominateur $4(4+h)$, comme dans la méthode
2 b).}{$f(4+h)=\frac{1}{4+h}$ et $f(4)=\frac{1}{4}$}
\begin{calculs}
f(4+h)-f(4) &= \dfrac{1}{4+h}-\dfrac{1}{4} &&\justif{on recopie la différence}\\
f(4+h)-f(4) &= \dfrac{1\fois{4}}{(4+h)\fois{4}}-\dfrac{1\fois{(4+h)}}{4\fois{(4+h)}} &&\justif{même dénominateur}\\
f(4+h)-f(4) &= \dfrac{4-4-h}{4(4+h)} &&\justif{le \og$-$\fg{} porte sur tout $4+h$}\\
f(4+h)-f(4) &= \dfrac{-h}{4(4+h)} &&\justif{c'est l'égalité demandée}
\end{calculs}
\explique{On divise par $h$, puis on fait tendre $h$ vers $0$.}{définition 2}
\begin{calculs}
T(h) &= \dfrac{-\barre{h}}{4(4+h)}\times\dfrac{1}{\barre{h}} &&\justif{diviser par $h$, c'est multiplier par $\frac{1}{h}$}\\
T(h) &= \dfrac{-1}{4(4+h)} &&\\
f'(4) &= \lim_{h\to 0}\dfrac{-1}{4(4+h)} &&\\
f'(4) &= \dfrac{-1}{4\times 4} &&\justif{$4+h$ tend vers $4$}\\
f'(4) &= \res{-\dfrac{1}{16}} &&
\end{calculs}
\cas{2.}
\begin{calculs}
g(1) &= \dfrac{1}{1+1} = \dfrac{1}{2} &&\justif{on remplace $x$ par $1$}\\
g(1+h) &= \dfrac{1}{(1+h)+1} &&\justif{on remplace $x$ par $1+h$}\\
g(1+h) &= \dfrac{1}{2+h} &&\\
g(1+h)-g(1) &= \dfrac{1}{2+h}-\dfrac{1}{2} &&\\
g(1+h)-g(1) &= \dfrac{1\fois{2}}{(2+h)\fois{2}}-\dfrac{1\fois{(2+h)}}{2\fois{(2+h)}} &&\justif{même dénominateur}\\
g(1+h)-g(1) &= \dfrac{2-2-h}{2(2+h)} &&\\
g(1+h)-g(1) &= \res{\dfrac{-h}{2(2+h)}} &&
\end{calculs}
\begin{calculs}
T(h) &= \dfrac{-\barre{h}}{2(2+h)}\times\dfrac{1}{\barre{h}} &&\justif{on divise par $h$}\\
T(h) &= \dfrac{-1}{2(2+h)} &&\\
g'(1) &= \lim_{h\to 0}\dfrac{-1}{2(2+h)} &&\\
g'(1) &= \dfrac{-1}{2\times 2} &&\justif{$2+h$ tend vers $2$}\\
g'(1) &= \res{-\dfrac{1}{4}} &&
\end{calculs}
\rem{On ne soustrait jamais deux fractions en soustrayant leurs dénominateurs : on les met
d'abord au même dénominateur. Laisser le dénominateur factorisé, $2(2+h)$, permet de
simplifier par $h$ sans rien développer.}
\end{exo}

\begin{exo}[Pour un nombre $a$ quelconque]
\cas{1.}
\explique{Même calcul qu'à l'exercice 5, avec la lettre $a$ à la place d'un nombre : $a$
est fixé, seul $h$ varie.}{$a$ se traite comme un nombre connu}
\begin{calculs}
f(a) &= a^{2}-4a &&\justif{on remplace $x$ par $a$}\\
f(a+h) &= (a+h)^{2}-4(a+h) &&\justif{on remplace $x$ par $a+h$}\\
f(a+h) &= a^{2}+2ah+h^{2}-4a-4h &&\\
f(a+h)-f(a) &= a^{2}+2ah+h^{2}-4a-4h-(a^{2}-4a) &&\justif{parenthèses : on soustrait tout $f(a)$}\\
f(a+h)-f(a) &= a^{2}+2ah+h^{2}-4a-4h-a^{2}+4a &&\\
f(a+h)-f(a) &= 2ah+h^{2}-4h &&\justif{$a^{2}$ et $4a$ s'annulent}\\
T(h) &= \dfrac{\barre{h}(2a+h-4)}{\barre{h}} &&\justif{on factorise par $h$}\\
T(h) &= \res{2a-4+h} &&\justif{c'est l'expression demandée}
\end{calculs}
\cas{2.}
\begin{calculs}
f'(a) &= \lim_{h\to 0}\,(2a-4+h) &&\justif{limite de $T(h)$}\\
f'(a) &= \res{2a-4} &&\justif{$h$ tend vers $0$ ; $a$ ne bouge pas}
\end{calculs}
\cas{3.}
\begin{calculs}
f'(0) &= 2\times 0-4 = \res{-4} &&\justif{on remplace $a$ par $0$}\\
f'(2) &= 2\times 2-4 = \res{0} &&\\
f'(5) &= 2\times 5-4 = \res{6} &&
\end{calculs}
\rem{L'intérêt du calcul avec $a$ : une seule limite donne d'un coup tous les nombres
dérivés de $f$ ; on remplace ensuite $a$ par le nombre voulu.}
\end{exo}

\partie{Tangente à une courbe}

\begin{exo}[Lecture graphique]
\begin{center}
\begin{tikzpicture}[x=0.7cm,y=0.7cm]
  \repere{-4}{-4}{6}{4}
  \gradx{-3,-2,-1,1,2,3,4,5} \grady{-3,-2,-1,1,2,3}
  \draw[coulCourbe,line width=1.1pt]
    (-3,-2) .. controls (-2.333,1) and (-1.667,2) .. (-1,2)
            .. controls (-0.333,2) and (0.333,1) .. (1,0)
            .. controls (1.667,-1) and (2.333,-2) .. (3,-2)
            .. controls (3.667,-2) and (4.333,-1) .. (5,2);
  \draw[coulExo,line width=0.9pt] (-2.6,2) -- (0.6,2);
  \draw[coulExo,line width=0.9pt] (-1.4,3.6) -- (2.6,-2.4);
  \draw[coulExo,line width=0.9pt] (1.4,-2) -- (4.6,-2);
  \draw[coulRep,dashed,thin] (-1,0) -- (-1,2);
  \draw[coulRep,dashed,thin] (3,0) -- (3,-2) (0,-2) -- (1.4,-2);
  \draw[coulRep,line width=1pt,-{Stealth[length=2mm]}] (-1,3) --
    node[pos=0.3,above,font=\footnotesize]{$+2$} (1,3);
  \draw[coulRep,line width=1pt,-{Stealth[length=2mm]}] (1,3) --
    node[right,font=\footnotesize]{$-3$} (1,0.1);
  \path[coulCourbe] (-1,2) node[croix]{} (1,0) node[croix]{} (3,-2) node[croix]{};
  \path[coulRep] (-1,3) node[croix]{};
  \node[font=\footnotesize] at (-1.5,2.5) {$A$};
  \node[font=\footnotesize,right] at (1.1,-0.4) {$B$};
  \node[font=\footnotesize] at (3.5,-2.5) {$C$};
  \node[coulCourbe,font=\footnotesize] at (5.5,1.2) {$\Cf$};
\end{tikzpicture}
\end{center}
\cas{1.}
\explique{$f(-1)$ est l'ordonnée du point de $\Cf$ d'abscisse $-1$ : c'est $A(-1\,;2)$,
donc $f(-1)=\res{2}$.}{on part de $-1$ sur l'axe des abscisses, on monte jusqu'à la
courbe (pointillés)}
\explique{$B(1\,;0)$ est sur l'axe des abscisses : $f(1)=\res{0}$ ; $C(3\,;-2)$ donne
$f(3)=\res{-2}$.}{une image se lit sur la courbe}
\cas{2.}
\explique{Les tangentes en $A$ et en $C$ sont horizontales : $f'(-1)=\res{0}$ et
$f'(3)=\res{0}$.}{une droite horizontale a un coefficient directeur nul}
\explique{La tangente en $B$ passe par le point $(-1\,;3)$ et par $B(1\,;0)$ : quand on
avance de $2$, elle descend de $3$ (flèches).}{on choisit deux points de la tangente
sur des nœuds du quadrillage}
\begin{calculs}
f'(1) &= \dfrac{\text{montée}}{\text{avancée}} &&\justif{coefficient directeur de la tangente en $B$}\\
f'(1) &= \dfrac{-3}{2} &&\justif{une descente compte en négatif}\\
f'(1) &= \res{-1{,}5} &&
\end{calculs}
\rem{On lit $f(a)$ sur la courbe (une ordonnée) et $f'(a)$ sur la tangente (une pente) :
ce sont deux lectures différentes. Ici, $f(1)=0$, mais $f'(1)\ne 0$.}
\cas{3.}
\begin{calculs}
y &= f'(1)(x-1)+f(1) &&\justif{propriété 1, avec $a=1$}\\
y &= -\dfrac{3}{2}(x-1)+0 &&\justif{$f'(1)=-\frac{3}{2}$ et $f(1)=0$}\\
y &= \res{-\dfrac{3}{2}x+\dfrac{3}{2}} &&\justif{on développe}
\end{calculs}
\explique{Vérification : pour $x=-1$, $y=\frac{3}{2}+\frac{3}{2}=3$ ; on retrouve le
point $(-1\,;3)$ lu sur la figure.}{on contrôle avec un point de la tangente}
\end{exo}

\begin{exo}[Équation de la tangente]
\explique{Dans chaque cas, on recopie la formule $y=f'(a)(x-a)+f(a)$, on remplace, puis
on développe.}{propriété 1}
\cas{a)}
\begin{calculs}
y &= f'(2)(x-2)+f(2) &&\justif{on recopie la formule, avec $a=2$}\\
y &= 4(x-2)+3 &&\justif{$f'(2)=4$ et $f(2)=3$}\\
y &= 4x-8+3 &&\justif{on développe}\\
y &= \res{4x-5} &&
\end{calculs}
\cas{b)}
\begin{calculs}
y &= f'(-1)(x+1)+f(-1) &&\justif{$x-(-1)=x+1$}\\
y &= -3(x+1)+2 &&\\
y &= -3x-3+2 &&\\
y &= \res{-3x-1} &&
\end{calculs}
\cas{c)}
\begin{calculs}
y &= f'(0)(x-0)+f(0) &&\justif{$x-0=x$}\\
y &= \res{\dfrac{1}{2}x+5} &&
\end{calculs}
\cas{d)}
\begin{calculs}
y &= f'(4)(x-4)+f(4) &&\\
y &= 0\times(x-4)-1 &&\justif{$f'(4)=0$}\\
y &= \res{-1} &&\justif{tangente horizontale}
\end{calculs}
\rem{Ne pas inverser les rôles : $f'(a)$, la pente, multiplie $(x-a)$ ; $f(a)$,
l'ordonnée du point, s'ajoute. Contrôle : pour $x=a$, on doit retrouver $y=f(a)$ ; en a),
$4\times 2-5=3$.}
\end{exo}

\begin{exo}[Nombre dérivé, puis tangente]
\cas{1.}
\begin{calculs}
f(1) &= 1+1-2 &&\justif{on calcule d'abord $f(1)$}\\
f(1) &= 0 &&\\
f(1+h) &= (1+h)^{2}+(1+h)-2 &&\justif{on remplace $x$ par $1+h$}\\
f(1+h) &= 1+2h+h^{2}+1+h-2 &&\\
f(1+h) &= h^{2}+3h &&\\
T(h) &= \dfrac{h^{2}+3h-0}{h} &&\justif{$T(h)=\frac{f(1+h)-f(1)}{h}$}\\
T(h) &= \dfrac{\barre{h}(h+3)}{\barre{h}} &&\\
T(h) &= h+3 &&\\
f'(1) &= \lim_{h\to 0}\,(h+3) &&\\
f'(1) &= \res{3} &&
\end{calculs}
\explique{Tangente au point d'abscisse $1$ :}{propriété 1}
\begin{calculs}
y &= f'(1)(x-1)+f(1) &&\justif{on recopie la formule}\\
y &= 3(x-1)+0 &&\\
y &= \res{3x-3} &&
\end{calculs}
\cas{2.}
\begin{calculs}
g(-1) &= \dfrac{1}{-1} = -1 &&\\
g(-1+h)-g(-1) &= \dfrac{1}{-1+h}-(-1) &&\justif{on soustrait $g(-1)=-1$}\\
g(-1+h)-g(-1) &= \dfrac{1}{-1+h}+\dfrac{1\fois{(-1+h)}}{1\fois{(-1+h)}} &&\justif{même dénominateur}\\
g(-1+h)-g(-1) &= \dfrac{1-1+h}{-1+h} &&\\
g(-1+h)-g(-1) &= \dfrac{h}{-1+h} &&\\
T(h) &= \dfrac{\barre{h}}{-1+h}\times\dfrac{1}{\barre{h}} &&\justif{diviser par $h$, c'est multiplier par $\frac{1}{h}$}\\
T(h) &= \dfrac{1}{-1+h} &&\\
g'(-1) &= \lim_{h\to 0}\dfrac{1}{-1+h} &&\\
g'(-1) &= \res{-1} &&\justif{$-1+h$ tend vers $-1$}
\end{calculs}
\explique{Tangente au point d'abscisse $-1$ :}{}
\begin{calculs}
y &= g'(-1)(x+1)+g(-1) &&\justif{$x-(-1)=x+1$}\\
y &= -1\times(x+1)+(-1) &&\\
y &= -x-1-1 &&\\
y &= \res{-x-2} &&
\end{calculs}
\rem{Contrôle : $g'(-1)=-1<0$, la tangente descend, comme la courbe de la fonction
inverse, décroissante sur $\intoo{-\infty}{0}$.}
\end{exo}

\begin{exo}[Tangente passant par un point]
\explique{La tangente a pour coefficient directeur le nombre dérivé (définition 3). Comme
elle passe par deux points connus, on calcule son coefficient directeur avec ces deux
points.}{$\dfrac{y_{B}-y_{A}}{x_{B}-x_{A}}$}
\cas{1.}
\begin{calculs}
f'(1) &= \dfrac{y_{B}-y_{A}}{x_{B}-x_{A}} &&\justif{coefficient directeur de la droite $(AB)$}\\
f'(1) &= \dfrac{4-(-2)}{3-1} &&\\
f'(1) &= \dfrac{6}{2} = \res{3} &&
\end{calculs}
\begin{calculs}
y &= f'(1)(x-1)+f(1) &&\justif{$A(1\,;-2)$ est sur $\Cf$ : $f(1)=-2$}\\
y &= 3(x-1)-2 &&\\
y &= 3x-3-2 &&\\
y &= \res{3x-5} &&
\end{calculs}
\explique{Vérification : pour $x=3$, $y=9-5=4$ ; la droite passe bien par $B(3\,;4)$.}{}
\cas{2.}
\begin{calculs}
g'(-2) &= \dfrac{y_{D}-y_{C}}{x_{D}-x_{C}} &&\justif{coefficient directeur de $(CD)$}\\
g'(-2) &= \dfrac{-3-1}{0-(-2)} &&\\
g'(-2) &= \dfrac{-4}{2} = \res{-2} &&
\end{calculs}
\begin{calculs}
y &= g'(-2)(x+2)+g(-2) &&\justif{$g(-2)=1$ ; $x-(-2)=x+2$}\\
y &= -2(x+2)+1 &&\\
y &= -2x-4+1 &&\\
y &= \res{-2x-3} &&
\end{calculs}
\explique{Vérification : pour $x=0$, $y=-3$ ; c'est bien $D(0\,;-3)$.}{l'ordonnée à
l'origine se lit directement}
\rem{Le point $B$ est sur la tangente, pas sur la courbe : seul $A$ est commun aux deux.
On n'écrit donc pas $f(3)=4$.}
\end{exo}

\partie{Approximation linéaire}

\begin{exo}[Valeurs approchées sans calculatrice]
\cas{1.}
\explique{Nombre dérivé en $6$, par le taux d'accroissement :}{méthode 2}
\begin{calculs}
f(6) &= 36 &&\\
f(6+h) &= 36+12h+h^{2} &&\justif{$(a+b)^{2}=a^{2}+2ab+b^{2}$}\\
T(h) &= \dfrac{12h+h^{2}}{h} &&\justif{$T(h)=\frac{f(6+h)-f(6)}{h}$}\\
T(h) &= 12+h &&\justif{on simplifie par $h$}\\
f'(6) &= \res{12} &&\justif{limite quand $h$ tend vers $0$}
\end{calculs}
\explique{Valeur approchée, avec $a=6$ et $h=0{,}02$ :}{$a$ proche de $6{,}02$, où $f(a)$ se
calcule de tête}
\begin{calculs}
f(6+0{,}02) &\approx f(6)+f'(6)\times 0{,}02 &&\justif{on recopie la formule}\\
6{,}02^{2} &\approx 36+12\times 0{,}02 &&\\
6{,}02^{2} &\approx \res{36{,}24} &&\justif{$12\times 0{,}02=0{,}24$}
\end{calculs}
\cas{2.}
\explique{Nombre dérivé en $1$ :}{$(1+h)^{3}=(1+h)^{2}(1+h)$}
\begin{calculs}
g(1) &= 1 &&\\
g(1+h) &= (1+2h+h^{2})(1+h) &&\\
g(1+h) &= 1+h+2h+2h^{2}+h^{2}+h^{3} &&\justif{on développe}\\
g(1+h) &= 1+3h+3h^{2}+h^{3} &&\\
T(h) &= \dfrac{3h+3h^{2}+h^{3}}{h} &&\\
T(h) &= 3+3h+h^{2} &&\justif{on simplifie par $h$}\\
g'(1) &= \res{3} &&\justif{limite quand $h$ tend vers $0$}
\end{calculs}
\begin{calculs}
g(1+0{,}01) &\approx g(1)+g'(1)\times 0{,}01 &&\justif{$a=1$ et $h=0{,}01$}\\
1{,}01^{3} &\approx 1+3\times 0{,}01 &&\\
1{,}01^{3} &\approx \res{1{,}03} &&
\end{calculs}
\cas{3.}
\explique{On cherche un nombre $a$ proche de $0{,}98$ dont l'inverse se calcule de tête :
$a=1$, et $0{,}98=1+h$ avec $h=-0{,}02$.}{$h$ peut être négatif : on recule}
\begin{calculs}
k(1+h)-k(1) &= \dfrac{1}{1+h}-1 &&\justif{$k(1)=1$}\\
k(1+h)-k(1) &= \dfrac{1}{1+h}-\dfrac{1\fois{(1+h)}}{1\fois{(1+h)}} &&\justif{même dénominateur}\\
k(1+h)-k(1) &= \dfrac{-h}{1+h} &&\\
T(h) &= \dfrac{-\barre{h}}{1+h}\times\dfrac{1}{\barre{h}} &&\justif{on divise par $h$}\\
T(h) &= \dfrac{-1}{1+h} &&\\
k'(1) &= \res{-1} &&\justif{$1+h$ tend vers $1$}
\end{calculs}
\begin{calculs}
k(1-0{,}02) &\approx k(1)+k'(1)\times(-0{,}02) &&\justif{on recopie la formule}\\
\dfrac{1}{0{,}98} &\approx 1+(-1)\times(-0{,}02) &&\\
\dfrac{1}{0{,}98} &\approx \res{1{,}02} &&\justif{$(-1)\times(-0{,}02)=+0{,}02$}
\end{calculs}
\cas{4.}
\explique{On prend $a=2$, où $k(2)=0{,}5$, et $h=0{,}01$.}{même calcul que la méthode 2 b)
du cours}
\begin{calculs}
k(2+h)-k(2) &= \dfrac{1}{2+h}-\dfrac{1}{2} &&\\
k(2+h)-k(2) &= \dfrac{1\fois{2}}{(2+h)\fois{2}}-\dfrac{1\fois{(2+h)}}{2\fois{(2+h)}} &&\justif{même dénominateur}\\
k(2+h)-k(2) &= \dfrac{-h}{2(2+h)} &&\\
T(h) &= \dfrac{-1}{2(2+h)} &&\justif{on divise par $h$}\\
k'(2) &= \res{-\dfrac{1}{4}} &&\justif{$2+h$ tend vers $2$}
\end{calculs}
\begin{calculs}
k(2+0{,}01) &\approx k(2)+k'(2)\times 0{,}01 &&\\
\dfrac{1}{2{,}01} &\approx 0{,}5-\dfrac{1}{4}\times 0{,}01 &&\\
\dfrac{1}{2{,}01} &\approx 0{,}5-0{,}0025 &&\justif{$\frac{1}{4}\times 0{,}01=0{,}0025$}\\
\dfrac{1}{2{,}01} &\approx \res{0{,}4975} &&
\end{calculs}
\cas{5.}
\begin{center}
\renewcommand{\arraystretch}{2.4}
\begin{tabular}{|c|c|c|c|c|}
\hline
 & $6{,}02^{2}$ & $1{,}01^{3}$ & $\frac{1}{0{,}98}$ & $\frac{1}{2{,}01}$\\ \hline
valeur approchée & $36{,}24$ & $1{,}03$ & $1{,}02$ & $0{,}4975$\\ \hline
calculatrice & $36{,}2404$ & $1{,}030301$ & $1{,}020408\ldots$ & $0{,}497512\ldots$\\ \hline
\end{tabular}
\end{center}
\explique{Les quatre valeurs approchées sont \res{justes au centième} (et même mieux).}{les
écarts sont de l'ordre de $h^{2}$}
\rem{Pour $f(x)=x^{2}$, l'écart vaut exactement $h^{2}=0{,}0004$ : c'est le terme négligé
dans $f(6+h)=36+12h+h^{2}$. Le point clé : choisir $a$ proche, où $f(a)$ et $f'(a)$ se
calculent de tête, et $h$ petit (positif ou négatif).}
\end{exo}

\begin{exo}[Augmentations et baisses successives]
\cas{1.}
\explique{Augmenter de $1{,}5$~\% revient à multiplier par $1{,}015$ ; deux mois de suite,
par $1{,}015\times 1{,}015=1{,}015^{2}$.}{coefficient multiplicateur
$1+\frac{t}{100}$}
\begin{calculs}
f(1) &= 1 &&\\
f(1+h) &= 1+2h+h^{2} &&\justif{$(1+h)^{2}$}\\
T(h) &= \dfrac{2h+h^{2}}{h} &&\justif{$T(h)=\frac{f(1+h)-f(1)}{h}$}\\
T(h) &= 2+h &&\\
f'(1) &= \res{2} &&\justif{limite quand $h$ tend vers $0$}
\end{calculs}
\begin{calculs}
f(1+0{,}015) &\approx f(1)+f'(1)\times 0{,}015 &&\justif{$a=1$ et $h=0{,}015$}\\
1{,}015^{2} &\approx 1+2\times 0{,}015 &&\\
1{,}015^{2} &\approx \res{1{,}03} &&
\end{calculs}
\explique{Le coefficient multiplicateur sur deux mois vaut environ $1{,}03$ : l'abonnement
augmente d'environ $\res{3~\%}$.}{la calculatrice donne $1{,}030225$, soit $3{,}0225$~\%}
\cas{2.}
\explique{Baisser de $2$~\% revient à multiplier par $0{,}98=1-0{,}02$ ; trois ans de
suite, par $0{,}98^{3}$.}{$h=-0{,}02$ : on recule}
\begin{calculs}
g(1) &= 1 &&\\
T(h) &= 3+3h+h^{2} &&\justif{calcul de l'exercice 13, question 2}\\
g'(1) &= \res{3} &&\justif{limite quand $h$ tend vers $0$}
\end{calculs}
\begin{calculs}
g(1-0{,}02) &\approx g(1)+g'(1)\times(-0{,}02) &&\\
0{,}98^{3} &\approx 1+3\times(-0{,}02) &&\\
0{,}98^{3} &\approx \res{0{,}94} &&\justif{$1-0{,}06$}
\end{calculs}
\explique{Le coefficient multiplicateur vaut environ $0{,}94=1-0{,}06$ : la population
baisse d'environ $\res{6~\%}$ en $3$ ans.}{la calculatrice donne $0{,}941192$, soit une
baisse de $5{,}88$~\%}
\rem{Pour de petits pourcentages, plusieurs évolutions successives de $t~\%$ font à peu
près la somme des pourcentages : c'est l'approximation $1+2h$, $1+3h$ obtenue ici. Ce
n'est qu'une valeur approchée : pour de grands pourcentages, l'écart devient
important.}
\end{exo}

\begin{exo}[Courbe et tangente]
\cas{1.}
\begin{calculs}
f(2) &= 2^{2}+3\times 2 &&\\
f(2) &= \res{10} &&\\
f(2+h) &= (2+h)^{2}+3(2+h) &&\justif{on remplace $x$ par $2+h$}\\
f(2+h) &= 4+4h+h^{2}+6+3h &&\\
f(2+h) &= h^{2}+7h+10 &&\\
T(h) &= \dfrac{h^{2}+7h+10-10}{h} &&\justif{$T(h)=\frac{f(2+h)-f(2)}{h}$}\\
T(h) &= \dfrac{\barre{h}(h+7)}{\barre{h}} &&\\
T(h) &= h+7 &&\\
f'(2) &= \lim_{h\to 0}\,(h+7) &&\\
f'(2) &= \res{7} &&
\end{calculs}
\cas{2.}
\explique{$2{,}1=2+0{,}1$ : on prend $a=2$ et $h=0{,}1$.}{propriété 2}
\begin{calculs}
f(2+0{,}1) &\approx f(2)+f'(2)\times 0{,}1 &&\justif{on recopie la formule}\\
f(2{,}1) &\approx 10+7\times 0{,}1 &&\\
f(2{,}1) &\approx \res{10{,}7} &&
\end{calculs}
\begin{calculs}
f(2{,}1) &= 2{,}1^{2}+3\times 2{,}1 &&\justif{valeur exacte}\\
f(2{,}1) &= 4{,}41+6{,}3 &&\\
f(2{,}1) &= 10{,}71 &&
\end{calculs}
\explique{L'écart vaut $10{,}71-10{,}7=\res{0{,}01}$.}{c'est $h^{2}=0{,}1^{2}$, le terme
négligé dans $f(2+h)=10+7h+h^{2}$}
\cas{3.}
\explique{$1{,}95=2-0{,}05$ : on prend $a=2$ et $h=-0{,}05$.}{on recule}
\begin{calculs}
f(2-0{,}05) &\approx f(2)+f'(2)\times(-0{,}05) &&\\
f(1{,}95) &\approx 10-0{,}35 &&\justif{$7\times 0{,}05=0{,}35$}\\
f(1{,}95) &\approx \res{9{,}65} &&
\end{calculs}
\begin{calculs}
f(1{,}95) &= 1{,}95^{2}+3\times 1{,}95 &&\justif{valeur exacte}\\
f(1{,}95) &= 3{,}8025+5{,}85 &&\\
f(1{,}95) &= 9{,}6525 &&
\end{calculs}
\explique{L'écart vaut $9{,}6525-9{,}65=\res{0{,}0025}$.}{c'est $0{,}05^{2}$}
\rem{Plus $h$ est petit, meilleure est l'approximation : ici, diviser $h$ par $2$ divise
l'écart par $4$.}
\end{exo}

\partie{Dérivées des fonctions polynômes}

\begin{exo}[Dériver un polynôme]
\explique{On dérive terme à terme avec le tableau de la propriété 4 ; un coefficient se
garde.}{la dérivée de $kx^{n}$ est $k\times nx^{n-1}$}
\cas{a)}
\begin{calculs}
f'(x) &= \res{0} &&\justif{la dérivée d'une constante est nulle}
\end{calculs}
\explique{La courbe de $f$ est la droite horizontale d'équation $y=7$ : sa pente est nulle
partout.}{}
\cas{b)}
\begin{calculs}
g'(x) &= \res{5} &&\justif{$mx+p$ donne $m$ : la constante $-3$ disparaît}
\end{calculs}
\cas{c)}
\begin{calculs}
h'(x) &= \res{6x^{5}} &&\justif{$x^{n}$ donne $nx^{n-1}$, avec $n=6$}
\end{calculs}
\cas{d)}
\begin{calculs}
k'(x) &= 3\times 2x-4 &&\justif{$x^{2}$ donne $2x$ ; $-4x$ donne $-4$ ; $1$ donne $0$}\\
k'(x) &= \res{6x-4} &&
\end{calculs}
\cas{e)}
\begin{calculs}
m'(x) &= -2\times 3x^{2}+2x-5 &&\justif{$x^{3}$ donne $3x^{2}$ ; $x^{2}$ donne $2x$}\\
m'(x) &= \res{-6x^{2}+2x-5} &&
\end{calculs}
\cas{f)}
\begin{calculs}
p'(x) &= 4x^{3}-2\times 3x^{2}+7 &&\justif{$-1$ donne $0$}\\
p'(x) &= \res{4x^{3}-6x^{2}+7} &&
\end{calculs}
\rem{Le coefficient se garde : la dérivée de $-2x^{3}$ est $-2\times 3x^{2}=-6x^{2}$, ni
$-2x^{2}$ ni $3x^{2}$. L'exposant descend multiplier, puis baisse de $1$.}
\end{exo}

\begin{exo}[Avec des fractions]
\cas{a)}
\begin{calculs}
f'(x) &= \dfrac{1}{2}\times 2x-3 &&\justif{$\frac{1}{2}\times 2=1$}\\
f'(x) &= \res{x-3} &&
\end{calculs}
\cas{b)}
\begin{calculs}
g'(x) &= \dfrac{1}{3}\times 3x^{2}-2x &&\justif{la constante $4$ donne $0$}\\
g'(x) &= \res{x^{2}-2x} &&
\end{calculs}
\cas{c)}
\begin{calculs}
h'(x) &= \dfrac{1}{4}\times 4x^{3}+\dfrac{2}{3}\times 3x^{2} &&\\
h'(x) &= \res{x^{3}+2x^{2}} &&\justif{$\frac{2}{3}\times 3=2$}
\end{calculs}
\cas{d)}
\explique{On écrit d'abord $k(x)$ comme un polynôme : diviser par $3$, c'est diviser
chaque terme par $3$.}{}
\begin{calculs}
k(x) &= \dfrac{x^{3}}{3}-\dfrac{6x}{3} &&\\
k(x) &= \dfrac{1}{3}x^{3}-2x &&\\
k'(x) &= \dfrac{1}{3}\times 3x^{2}-2 &&\\
k'(x) &= \res{x^{2}-2} &&
\end{calculs}
\rem{Les coefficients $\frac{1}{2}$, $\frac{1}{3}$, $\frac{1}{4}$ se simplifient avec
l'exposant qui descend : $\frac{1}{3}x^{3}$ a pour dérivée $x^{2}$.}
\end{exo}

\begin{exo}[Développer d'abord]
\explique{Le tableau de la propriété 4 ne donne la dérivée ni d'un carré ni d'un produit :
on développe d'abord, pour obtenir un polynôme.}{}
\cas{a)}
\begin{calculs}
f(x) &= (x+4)^{2} &&\justif{on recopie l'expression}\\
f(x) &= x^{2}+8x+16 &&\justif{$(a+b)^{2}=a^{2}+2ab+b^{2}$}\\
f'(x) &= \res{2x+8} &&
\end{calculs}
\cas{b)}
\begin{calculs}
g(x) &= (3x-1)^{2} &&\\
g(x) &= 9x^{2}-6x+1 &&\justif{$(3x)^{2}=9x^{2}$ ; $2\times 3x\times 1=6x$}\\
g'(x) &= \res{18x-6} &&
\end{calculs}
\cas{c)}
\begin{calculs}
h(x) &= x(x^{2}-5) &&\\
h(x) &= x^{3}-5x &&\\
h'(x) &= \res{3x^{2}-5} &&
\end{calculs}
\cas{d)}
\begin{calculs}
k(x) &= (x+2)(x-2) &&\\
k(x) &= x^{2}-4 &&\justif{$(a+b)(a-b)=a^{2}-b^{2}$}\\
k'(x) &= \res{2x} &&
\end{calculs}
\rem{Erreur fréquente : dériver $(3x-1)^{2}$ \og comme $x^{2}$\fg{} en $2(3x-1)=6x-2$ ;
le développement donne $18x-6$.}
\end{exo}

\begin{exo}[Entraînement : polynômes]
\cas{a)}
\begin{calculs}
f'(x) &= 3x^{2}-7\times 2x+2 &&\justif{la constante $-9$ donne $0$}\\
f'(x) &= \res{3x^{2}-14x+2} &&
\end{calculs}
\cas{b)}
\begin{calculs}
g'(x) &= -4\times 5x^{4}+3\times 2x &&\\
g'(x) &= \res{-20x^{4}+6x} &&
\end{calculs}
\cas{c)}
\begin{calculs}
h'(x) &= 6\times 4x^{3}-3x^{2}+1 &&\\
h'(x) &= \res{24x^{3}-3x^{2}+1} &&
\end{calculs}
\cas{d)}
\begin{calculs}
k'(x) &= \dfrac{2}{5}\times 5x^{4}-1 &&\justif{$\frac{2}{5}\times 5=2$}\\
k'(x) &= \res{2x^{4}-1} &&
\end{calculs}
\cas{e)}
\begin{calculs}
m'(x) &= 0{,}5\times 4x^{3}-1{,}5\times 2x &&\\
m'(x) &= \res{2x^{3}-3x} &&
\end{calculs}
\cas{f)}
\begin{calculs}
p(x) &= \dfrac{1}{4}x^{2}-\dfrac{3}{4} &&\justif{on divise chaque terme par $4$}\\
p'(x) &= \dfrac{1}{4}\times 2x &&\justif{$-\frac{3}{4}$ est une constante}\\
p'(x) &= \res{\dfrac{1}{2}x} &&
\end{calculs}
\cas{g)}
\begin{calculs}
q'(x) &= \res{7x^{6}-5x^{4}+3x^{2}} &&\justif{$x^{n}$ donne $nx^{n-1}$, terme à terme}
\end{calculs}
\cas{h)}
\begin{calculs}
r(x) &= 3x^{2}-6x+15 &&\justif{on développe}\\
r'(x) &= \res{6x-6} &&
\end{calculs}
\rem{En f) et h), un nombre qui multiplie toute l'expression se garde aussi : sans
développer, $r'(x)=3(2x-2)$.}
\end{exo}

\begin{exo}[Fonction dérivée et nombres dérivés]
\cas{1.}
\begin{calculs}
f'(x) &= 3x^{2}-3\times 2x &&\justif{la constante $2$ donne $0$}\\
f'(x) &= \res{3x^{2}-6x} &&
\end{calculs}
\cas{2.}
\explique{$f'(x)$ donne le nombre dérivé en tout $x$ : on remplace $x$ par le nombre
voulu, sans refaire de limite.}{définition 4}
\begin{calculs}
f'(0) &= 3\times 0^{2}-6\times 0 = \res{0} &&\\
f'(1) &= 3\times 1^{2}-6\times 1 = \res{-3} &&\\
f'(-1) &= 3\times(-1)^{2}-6\times(-1) = \res{9} &&\justif{$(-1)^{2}=1$ et $-6\times(-1)=+6$}\\
f'(2) &= 3\times 2^{2}-6\times 2 = \res{0} &&\justif{$12-12$}
\end{calculs}
\cas{3.}
\begin{calculs}
f(1) &= 1-3+2 = 0 &&\justif{on calcule d'abord $f(1)$}\\
y &= f'(1)(x-1)+f(1) &&\justif{propriété 1}\\
y &= -3(x-1)+0 &&\justif{$f'(1)=-3$ (question 2)}\\
y &= \res{-3x+3} &&
\end{calculs}
\cas{4.}
\explique{Une tangente est horizontale quand son coefficient directeur est nul : on résout
$f'(x)=0$.}{le coefficient directeur de la tangente est $f'(a)$}
\begin{calculs}
3x^{2}-6x &= 0 &&\justif{on recopie l'équation}\\
3x(x-2) &= 0 &&\justif{on factorise par $3x$}
\end{calculs}
\explique{Un produit est nul si l'un de ses facteurs est nul : $x=0$ ou $x=2$.}{on le
voyait déjà à la question 2}
\begin{calculs}
f(0) &= 2 &&\\
f(2) &= 8-12+2 = -2 &&
\end{calculs}
\explique{La tangente est horizontale aux points $\res{(0\,;2)}$ et
$\res{(2\,;-2)}$.}{on demande des points : on donne les deux coordonnées}
\end{exo}

\partie{Dérivées des autres fonctions de référence}

\begin{exo}[Nombres dérivés des fonctions de référence]
\explique{On écrit d'abord la fonction dérivée (propriété 5), puis on remplace
$x$.}{méthode 7}
\cas{a)}
\begin{calculs}
f'(x) &= -\dfrac{1}{x^{2}} &&\justif{on écrit d'abord la dérivée}\\
f'(5) &= -\dfrac{1}{5^{2}} = \res{-\dfrac{1}{25}} &&\justif{puis on remplace $x$}\\
f'\Bigl(\dfrac{1}{2}\Bigr) &= -\dfrac{1}{\left(\frac{1}{2}\right)^{2}} &&\\
f'\Bigl(\dfrac{1}{2}\Bigr) &= -\dfrac{1}{\frac{1}{4}} &&\justif{$\frac{1}{2}\times\frac{1}{2}=\frac{1}{4}$}\\
f'\Bigl(\dfrac{1}{2}\Bigr) &= \res{-4} &&\justif{diviser par $\frac{1}{4}$, c'est multiplier par $4$}
\end{calculs}
\cas{b)}
\begin{calculs}
g'(x) &= \dfrac{1}{2\sqrt{x}} &&\\
g'(16) &= \dfrac{1}{2\times 4} = \res{\dfrac{1}{8}} &&\justif{$\sqrt{16}=4$}\\
g'\Bigl(\dfrac{1}{4}\Bigr) &= \dfrac{1}{2\times\frac{1}{2}} &&\justif{$\sqrt{\frac{1}{4}}=\frac{1}{2}$, car $\frac{1}{2}\times\frac{1}{2}=\frac{1}{4}$}\\
g'\Bigl(\dfrac{1}{4}\Bigr) &= \res{1} &&
\end{calculs}
\cas{c)}
\begin{calculs}
k'(x) &= -\dfrac{2}{x^{3}} &&\justif{$\frac{1}{x^{n}}$ donne $-\frac{n}{x^{n+1}}$, avec $n=2$}\\
k'(1) &= -\dfrac{2}{1^{3}} = \res{-2} &&\\
k'(-2) &= -\dfrac{2}{(-2)^{3}} &&\\
k'(-2) &= -\dfrac{2}{-8} = \res{\dfrac{1}{4}} &&\justif{$(-2)^{3}=-8$ : un cube garde le signe}
\end{calculs}
\rem{Erreur fréquente : $\frac{1}{x^{2}}$ donne $-\frac{2}{x^{3}}$, ni $-\frac{2}{x}$ ni
$-\frac{1}{2x}$ : l'exposant du dénominateur augmente de $1$.}
\end{exo}

\begin{exo}[Dériver avec les fonctions de référence]
\cas{a)}
\begin{calculs}
f(x) &= 4\times\dfrac{1}{x} &&\justif{le coefficient $4$ se garde}\\
f'(x) &= 4\times\left(-\dfrac{1}{x^{2}}\right) &&\\
f'(x) &= \res{-\dfrac{4}{x^{2}}} &&
\end{calculs}
\cas{b)}
\begin{calculs}
g'(x) &= -3\times\dfrac{1}{2\sqrt{x}} &&\justif{$\sqrt{x}$ donne $\frac{1}{2\sqrt{x}}$}\\
g'(x) &= \res{-\dfrac{3}{2\sqrt{x}}} &&
\end{calculs}
\cas{c)}
\begin{calculs}
h'(x) &= \res{-\dfrac{4}{x^{5}}} &&\justif{$\frac{1}{x^{n}}$ donne $-\frac{n}{x^{n+1}}$, avec $n=4$}
\end{calculs}
\cas{d)}
\begin{calculs}
k(x) &= 5\times\dfrac{1}{x^{3}} &&\\
k'(x) &= 5\times\left(-\dfrac{3}{x^{4}}\right) &&\justif{$n=3$}\\
k'(x) &= \res{-\dfrac{15}{x^{4}}} &&
\end{calculs}
\cas{e)}
\begin{calculs}
m'(x) &= 2\times(-\sin x) &&\justif{$\cos x$ donne $-\sin x$}\\
m'(x) &= \res{-2\sin x} &&
\end{calculs}
\cas{f)}
\begin{calculs}
p'(x) &= -1\times\cos x &&\justif{$\sin x$ donne $\cos x$}\\
p'(x) &= \res{-\cos x} &&
\end{calculs}
\rem{Ne pas confondre : $\cos x$ donne $-\sin x$, mais $\sin x$ donne $+\cos x$. Le
\og$-$\fg{} n'apparaît qu'en dérivant le cosinus.}
\end{exo}

\begin{exo}[Entraînement : fonctions de référence]
\cas{a)}
\begin{calculs}
f(x) &= -7\times\dfrac{1}{x} &&\\
f'(x) &= -7\times\left(-\dfrac{1}{x^{2}}\right) &&\\
f'(x) &= \res{\dfrac{7}{x^{2}}} &&\justif{les deux signes \og$-$\fg{} donnent \og$+$\fg}
\end{calculs}
\cas{b)}
\begin{calculs}
g'(x) &= 8\times\dfrac{1}{2\sqrt{x}} &&\\
g'(x) &= \res{\dfrac{4}{\sqrt{x}}} &&\justif{on simplifie par $2$}
\end{calculs}
\cas{c)}
\begin{calculs}
h(x) &= 2\times\dfrac{1}{x^{2}} &&\\
h'(x) &= 2\times\left(-\dfrac{2}{x^{3}}\right) &&\\
h'(x) &= \res{-\dfrac{4}{x^{3}}} &&
\end{calculs}
\cas{d)}
\begin{calculs}
k(x) &= -1\times\dfrac{1}{x^{5}} &&\\
k'(x) &= -1\times\left(-\dfrac{5}{x^{6}}\right) &&\justif{$n=5$}\\
k'(x) &= \res{\dfrac{5}{x^{6}}} &&
\end{calculs}
\cas{e)}
\begin{calculs}
m'(x) &= \dfrac{1}{3}\times\dfrac{1}{2\sqrt{x}} &&\\
m'(x) &= \res{\dfrac{1}{6\sqrt{x}}} &&\justif{$3\times 2=6$}
\end{calculs}
\cas{f)}
\begin{calculs}
p'(x) &= -4\times(-\sin x) &&\\
p'(x) &= \res{4\sin x} &&
\end{calculs}
\cas{g)}
\explique{$q(x)=\dfrac{1}{2x}$ n'est pas $\dfrac{1}{2}x$ : on l'écrit
$\dfrac{1}{2}\times\dfrac{1}{x}$.}{$\frac{1}{2x}=\frac{1\times 1}{2\times x}$}
\begin{calculs}
q(x) &= \dfrac{1}{2}\times\dfrac{1}{x} &&\\
q'(x) &= \dfrac{1}{2}\times\left(-\dfrac{1}{x^{2}}\right) &&\\
q'(x) &= \res{-\dfrac{1}{2x^{2}}} &&
\end{calculs}
\cas{h)}
\begin{calculs}
r'(x) &= \res{\dfrac{1}{2}\cos x} &&\justif{le coefficient $\frac{1}{2}$ se garde}
\end{calculs}
\rem{Deux signes \og$-$\fg{} donnent un \og$+$\fg{} : c'est l'erreur la plus fréquente en
a), d) et f).}
\end{exo}

\begin{exo}[Tangentes aux courbes de référence]
\cas{1.}
\begin{calculs}
f(1) &= \dfrac{1}{1} = 1 &&\\
f'(x) &= -\dfrac{1}{x^{2}} &&\justif{on écrit d'abord la dérivée}\\
f'(1) &= -\dfrac{1}{1^{2}} = -1 &&\\
y &= f'(1)(x-1)+f(1) &&\justif{propriété 1}\\
y &= -1\times(x-1)+1 &&\\
y &= -x+1+1 &&\\
y &= \res{-x+2} &&
\end{calculs}
\cas{2.}
\begin{calculs}
g(4) &= \sqrt{4} = 2 &&\\
g'(x) &= \dfrac{1}{2\sqrt{x}} &&\\
g'(4) &= \dfrac{1}{2\times 2} = \dfrac{1}{4} &&\\
y &= g'(4)(x-4)+g(4) &&\\
y &= \dfrac{1}{4}(x-4)+2 &&\\
y &= \dfrac{1}{4}x-1+2 &&\justif{$\frac{1}{4}\times 4=1$}\\
y &= \res{\dfrac{1}{4}x+1} &&
\end{calculs}
\cas{3.}
\explique{Une tangente horizontale a un coefficient directeur nul : il faudrait un nombre
$a$ avec $f'(a)=0$.}{le coefficient directeur de la tangente est $f'(a)$}
\explique{Or, pour $x\ne 0$, $x^{2}>0$, donc $\dfrac{1}{x^{2}}>0$ et
$f'(x)=-\dfrac{1}{x^{2}}<0$.}{un carré non nul est strictement positif}
\explique{$f'(x)$ ne vaut jamais $0$ : \res{la courbe de $f$ n'a aucune tangente
horizontale}.}{toutes ses tangentes descendent}
\rem{Une fraction n'est nulle que si son numérateur est nul ; ici, le numérateur vaut
$-1$.}
\end{exo}

\partie{Dérivée d'une somme et d'un produit}

\begin{exo}[Dériver une somme]
\explique{Chaque fonction est une somme : on dérive chaque terme.}{$(u+v)'=u'+v'$}
\cas{a)}
\begin{calculs}
f'(x) &= 3x^{2}+\left(-\dfrac{1}{x^{2}}\right) &&\justif{$x^{3}$ donne $3x^{2}$ ; $\frac{1}{x}$ donne $-\frac{1}{x^{2}}$}\\
f'(x) &= \res{3x^{2}-\dfrac{1}{x^{2}}} &&
\end{calculs}
\cas{b)}
\begin{calculs}
g'(x) &= 2\times 2x-\dfrac{1}{2\sqrt{x}} &&\justif{$\sqrt{x}$ donne $\frac{1}{2\sqrt{x}}$}\\
g'(x) &= \res{4x-\dfrac{1}{2\sqrt{x}}} &&
\end{calculs}
\cas{c)}
\begin{calculs}
h'(x) &= 3\times\left(-\dfrac{1}{x^{2}}\right)+5 &&\justif{$\frac{3}{x}=3\times\frac{1}{x}$}\\
h'(x) &= \res{-\dfrac{3}{x^{2}}+5} &&
\end{calculs}
\cas{d)}
\begin{calculs}
k'(x) &= 1+(-\sin x) &&\justif{$x$ donne $1$ ; $\cos x$ donne $-\sin x$}\\
k'(x) &= \res{1-\sin x} &&
\end{calculs}
\rem{$f$ et $h$ se dérivent pour $x\ne 0$, $g$ pour $x>0$ : la racine carrée n'est
dérivable que sur $\intoo{0}{+\infty}$.}
\end{exo}

\begin{exo}[Dériver un produit]
\explique{On nomme les deux facteurs $u$ et $v$, on calcule $u'$ et $v'$, puis on
applique la formule.}{$(uv)'=u'v+uv'$}
\cas{a)}
\begin{calculs}
u(x) &= 3x-2 &&\justif{$u'(x)=3$}\\
v(x) &= x+4 &&\justif{$v'(x)=1$}\\
f'(x) &= 3(x+4)+(3x-2)\times 1 &&\justif{$u'v+uv'$}\\
f'(x) &= 3x+12+3x-2 &&\\
f'(x) &= \res{6x+10} &&
\end{calculs}
\explique{Vérification : en développant, $f(x)=3x^{2}+10x-8$, de dérivée
$6x+10$.}{même résultat}
\cas{b)}
\begin{calculs}
u(x) &= x^{2}+1 &&\justif{$u'(x)=2x$}\\
v(x) &= 2x-3 &&\justif{$v'(x)=2$}\\
g'(x) &= 2x(2x-3)+(x^{2}+1)\times 2 &&\justif{$u'v+uv'$}\\
g'(x) &= 4x^{2}-6x+2x^{2}+2 &&\\
g'(x) &= \res{6x^{2}-6x+2} &&
\end{calculs}
\cas{c)}
\begin{calculs}
u(x) &= x+1 &&\justif{$u'(x)=1$}\\
v(x) &= \sqrt{x} &&\justif{$v'(x)=\frac{1}{2\sqrt{x}}$, pour $x>0$}\\
h'(x) &= 1\times\sqrt{x}+(x+1)\times\dfrac{1}{2\sqrt{x}} &&\justif{$u'v+uv'$}\\
h'(x) &= \dfrac{\sqrt{x}\fois{2\sqrt{x}}}{1\fois{2\sqrt{x}}}+\dfrac{x+1}{2\sqrt{x}} &&\justif{même dénominateur $2\sqrt{x}$}\\
h'(x) &= \dfrac{2x}{2\sqrt{x}}+\dfrac{x+1}{2\sqrt{x}} &&\justif{$\sqrt{x}\times\sqrt{x}=x$}\\
h'(x) &= \res{\dfrac{3x+1}{2\sqrt{x}}} &&
\end{calculs}
\explique{La forme $\sqrt{x}+\dfrac{x+1}{2\sqrt{x}}$ est déjà juste ; le même
dénominateur donne une seule fraction, plus simple à utiliser.}{}
\cas{d)}
\begin{calculs}
u(x) &= x^{2} &&\justif{$u'(x)=2x$}\\
v(x) &= \cos x &&\justif{$v'(x)=-\sin x$}\\
k'(x) &= 2x\cos x+x^{2}\times(-\sin x) &&\justif{$u'v+uv'$}\\
k'(x) &= \res{2x\cos x-x^{2}\sin x} &&
\end{calculs}
\rem{Erreur fréquente : $(uv)'=u'v'$. En a), on obtiendrait $3\times 1=3$, alors que la
dérivée est $6x+10$.}
\end{exo}

\begin{exo}[Entraînement : sommes et produits]
\cas{a)}
\begin{calculs}
f'(x) &= 2x-3+4\times\left(-\dfrac{1}{x^{2}}\right) &&\justif{somme ; $\frac{4}{x}=4\times\frac{1}{x}$}\\
f'(x) &= \res{2x-3-\dfrac{4}{x^{2}}} &&
\end{calculs}
\cas{b)}
\begin{calculs}
g'(x) &= 4\times\dfrac{1}{2\sqrt{x}}-2\times 3x^{2} &&\justif{somme}\\
g'(x) &= \res{\dfrac{2}{\sqrt{x}}-6x^{2}} &&\justif{on simplifie par $2$}
\end{calculs}
\cas{c)}
\begin{calculs}
u(x) &= x^{2}-4 &&\justif{$u'(x)=2x$}\\
v(x) &= 3x+1 &&\justif{$v'(x)=3$}\\
h'(x) &= 2x(3x+1)+(x^{2}-4)\times 3 &&\justif{$u'v+uv'$}\\
h'(x) &= 6x^{2}+2x+3x^{2}-12 &&\\
h'(x) &= \res{9x^{2}+2x-12} &&
\end{calculs}
\cas{d)}
\begin{calculs}
u(x) &= 2x-1 &&\justif{$u'(x)=2$}\\
v(x) &= \sqrt{x} &&\justif{$v'(x)=\frac{1}{2\sqrt{x}}$}\\
k'(x) &= 2\sqrt{x}+(2x-1)\times\dfrac{1}{2\sqrt{x}} &&\justif{$u'v+uv'$}\\
k'(x) &= \dfrac{2\sqrt{x}\fois{2\sqrt{x}}}{1\fois{2\sqrt{x}}}+\dfrac{2x-1}{2\sqrt{x}} &&\justif{même dénominateur}\\
k'(x) &= \dfrac{4x+2x-1}{2\sqrt{x}} &&\justif{$2\sqrt{x}\times 2\sqrt{x}=4x$}\\
k'(x) &= \res{\dfrac{6x-1}{2\sqrt{x}}} &&
\end{calculs}
\cas{e)}
\begin{calculs}
u(x) &= x+2 &&\justif{$u'(x)=1$}\\
v(x) &= \cos x &&\justif{$v'(x)=-\sin x$}\\
m'(x) &= 1\times\cos x+(x+2)\times(-\sin x) &&\justif{$u'v+uv'$}\\
m'(x) &= \res{\cos x-(x+2)\sin x} &&
\end{calculs}
\cas{f)}
\begin{calculs}
u(x) &= \dfrac{1}{x} &&\justif{$u'(x)=-\frac{1}{x^{2}}$}\\
v(x) &= x+5 &&\justif{$v'(x)=1$}\\
p'(x) &= -\dfrac{1}{x^{2}}\times(x+5)+\dfrac{1}{x}\times 1 &&\justif{$u'v+uv'$}\\
p'(x) &= \dfrac{-x-5}{x^{2}}+\dfrac{1\fois{x}}{x\fois{x}} &&\justif{même dénominateur $x^{2}$}\\
p'(x) &= \dfrac{-x-5+x}{x^{2}} &&\\
p'(x) &= \res{-\dfrac{5}{x^{2}}} &&
\end{calculs}
\rem{Plus court en f) : $p(x)=\dfrac{x+5}{x}=1+\dfrac{5}{x}$, puis on dérive la somme.}
\cas{g)}
\begin{calculs}
u(x) &= x^{2}+x &&\justif{$u'(x)=2x+1$}\\
v(x) &= x^{2}-x &&\justif{$v'(x)=2x-1$}\\
q'(x) &= (2x+1)(x^{2}-x)+(x^{2}+x)(2x-1) &&\justif{$u'v+uv'$}\\
q'(x) &= (2x^{3}-x^{2}-x)+(2x^{3}+x^{2}-x) &&\justif{on développe chaque produit}\\
q'(x) &= \res{4x^{3}-2x} &&
\end{calculs}
\rem{Vérification : $q(x)=(x^{2}+x)(x^{2}-x)=x^{4}-x^{2}$ (identité $(a+b)(a-b)$, avec
$a=x^{2}$ et $b=x$), dont la dérivée est bien $4x^{3}-2x$.}
\cas{h)}
\begin{calculs}
u(x) &= \sin x &&\justif{$u'(x)=\cos x$}\\
v(x) &= \cos x &&\justif{$v'(x)=-\sin x$}\\
r'(x) &= \cos x\times\cos x+\sin x\times(-\sin x) &&\justif{$u'v+uv'$}\\
r'(x) &= \res{(\cos x)^{2}-(\sin x)^{2}} &&
\end{calculs}
\end{exo}

\begin{exo}[Deux façons de dériver]
\cas{1.}
\begin{calculs}
u(x) &= 2x+5 &&\justif{$u'(x)=2$}\\
v(x) &= 1-3x &&\justif{$v'(x)=-3$}\\
f'(x) &= 2(1-3x)+(2x+5)\times(-3) &&\justif{$u'v+uv'$}\\
f'(x) &= 2-6x-6x-15 &&\\
f'(x) &= \res{-12x-13} &&
\end{calculs}
\cas{2.}
\begin{calculs}
f(x) &= (2x+5)(1-3x) &&\justif{on recopie l'expression}\\
f(x) &= 2x-6x^{2}+5-15x &&\justif{double distributivité}\\
f(x) &= -6x^{2}-13x+5 &&\\
f'(x) &= -6\times 2x-13 &&\\
f'(x) &= \res{-12x-13} &&
\end{calculs}
\cas{3.}
\explique{Les deux méthodes donnent la même dérivée, $\res{f'(x)=-12x-13}$.}{c'est un bon
moyen de vérifier une dérivée de produit}
\rem{Pour un produit de deux polynômes, développer d'abord est souvent aussi rapide ; la
formule $u'v+uv'$ devient indispensable avec $\sqrt{x}$, $\cos x$ ou $\sin x$, qu'on ne
sait pas développer.}
\end{exo}

\begin{exo}[Produit et tangente]
\cas{1.}
\begin{calculs}
u(x) &= x-3 &&\justif{$u'(x)=1$}\\
v(x) &= \sqrt{x} &&\justif{$v'(x)=\frac{1}{2\sqrt{x}}$}\\
f'(x) &= 1\times\sqrt{x}+(x-3)\times\dfrac{1}{2\sqrt{x}} &&\justif{$u'v+uv'$}\\
f'(x) &= \dfrac{\sqrt{x}\fois{2\sqrt{x}}}{1\fois{2\sqrt{x}}}+\dfrac{x-3}{2\sqrt{x}} &&\justif{même dénominateur}\\
f'(x) &= \dfrac{2x+x-3}{2\sqrt{x}} &&\justif{$\sqrt{x}\times\sqrt{x}=x$}\\
f'(x) &= \res{\dfrac{3x-3}{2\sqrt{x}}} &&\justif{c'est l'expression demandée}
\end{calculs}
\cas{2.}
\begin{calculs}
f(4) &= (4-3)\sqrt{4} &&\\
f(4) &= 1\times 2 = \res{2} &&\\
f'(4) &= \dfrac{3\times 4-3}{2\sqrt{4}} &&\justif{on utilise la forme de la question 1}\\
f'(4) &= \dfrac{9}{2\times 2} = \res{\dfrac{9}{4}} &&
\end{calculs}
\cas{3.}
\begin{calculs}
y &= f'(4)(x-4)+f(4) &&\justif{propriété 1}\\
y &= \dfrac{9}{4}(x-4)+2 &&\\
y &= \dfrac{9}{4}x-9+2 &&\justif{$\frac{9}{4}\times 4=9$}\\
y &= \res{\dfrac{9}{4}x-7} &&
\end{calculs}
\explique{Vérification : pour $x=4$, $y=9-7=2$ ; on retrouve $f(4)$.}{la tangente passe
par le point de la courbe}
\end{exo}

\partie{Dérivée d'un inverse et d'un quotient}

\begin{exo}[Dériver un inverse]
\explique{Chaque fonction est l'inverse d'une fonction $v$ : on nomme $v$, on calcule $v'$,
puis on applique la formule.}{$\left(\dfrac{1}{v}\right)'=-\dfrac{v'}{v^{2}}$}
\cas{a)}
\begin{calculs}
v(x) &= 2x+3 &&\justif{$v'(x)=2$}\\
f'(x) &= \res{-\dfrac{2}{(2x+3)^{2}}} &&\justif{$-\frac{v'}{v^{2}}$ : le dénominateur est élevé au carré}
\end{calculs}
\explique{$f$ est définie et dérivable pour $2x+3\ne 0$, c'est-à-dire
$x\ne-\dfrac{3}{2}$.}{là où $v$ ne s'annule pas}
\cas{b)}
\begin{calculs}
v(x) &= x^{2}+4 &&\justif{$v'(x)=2x$ ; $x^{2}+4>0$ : défini sur $\R$}\\
g'(x) &= \res{-\dfrac{2x}{(x^{2}+4)^{2}}} &&
\end{calculs}
\cas{c)}
\explique{$h(x)=\dfrac{3}{x-1}=3\times\dfrac{1}{x-1}$ : le coefficient $3$ se
garde.}{$(ku)'=ku'$}
\begin{calculs}
v(x) &= x-1 &&\justif{$v'(x)=1$}\\
h'(x) &= 3\times\left(-\dfrac{1}{(x-1)^{2}}\right) &&\\
h'(x) &= \res{-\dfrac{3}{(x-1)^{2}}} &&
\end{calculs}
\cas{d)}
\begin{calculs}
v(x) &= x^{2}-2x &&\justif{$v'(x)=2x-2$}\\
k'(x) &= \res{-\dfrac{2x-2}{(x^{2}-2x)^{2}}} &&\justif{tout $v'$ passe au numérateur}
\end{calculs}
\rem{Erreur fréquente : oublier le carré au dénominateur, ou écrire $-\dfrac{1}{v'}$. Le
\og$-$\fg{} et le carré viennent de $\left(\dfrac{1}{x}\right)'=-\dfrac{1}{x^{2}}$.}
\end{exo}

\begin{exo}[Dériver un quotient]
\explique{On nomme le numérateur $u$ et le dénominateur $v$, on calcule $u'$ et $v'$, puis
on applique la formule.}{$\left(\dfrac{u}{v}\right)'=\dfrac{u'v-uv'}{v^{2}}$}
\cas{a)}
\begin{calculs}
u(x) &= x+3 &&\justif{$u'(x)=1$}\\
v(x) &= x-2 &&\justif{$v'(x)=1$}\\
f'(x) &= \dfrac{1\times(x-2)-(x+3)\times 1}{(x-2)^{2}} &&\justif{$\frac{u'v-uv'}{v^{2}}$ : le dénominateur reste au carré}\\
f'(x) &= \dfrac{x-2-x-3}{(x-2)^{2}} &&\justif{le \og$-$\fg{} change les signes de $uv'$}\\
f'(x) &= \res{\dfrac{-5}{(x-2)^{2}}} &&
\end{calculs}
\cas{b)}
\begin{calculs}
u(x) &= 4x-1 &&\justif{$u'(x)=4$}\\
v(x) &= 2x+1 &&\justif{$v'(x)=2$}\\
g'(x) &= \dfrac{4(2x+1)-(4x-1)\times 2}{(2x+1)^{2}} &&\justif{$\frac{u'v-uv'}{v^{2}}$}\\
g'(x) &= \dfrac{8x+4-8x+2}{(2x+1)^{2}} &&\justif{$-(4x-1)\times 2=-8x+2$}\\
g'(x) &= \res{\dfrac{6}{(2x+1)^{2}}} &&
\end{calculs}
\cas{c)}
\begin{calculs}
u(x) &= x^{2} &&\justif{$u'(x)=2x$}\\
v(x) &= x-3 &&\justif{$v'(x)=1$}\\
h'(x) &= \dfrac{2x(x-3)-x^{2}\times 1}{(x-3)^{2}} &&\\
h'(x) &= \dfrac{2x^{2}-6x-x^{2}}{(x-3)^{2}} &&\\
h'(x) &= \res{\dfrac{x^{2}-6x}{(x-3)^{2}}} &&
\end{calculs}
\cas{d)}
\begin{calculs}
u(x) &= x^{2}-1 &&\justif{$u'(x)=2x$}\\
v(x) &= x^{2}+1 &&\justif{$v'(x)=2x$}\\
k'(x) &= \dfrac{2x(x^{2}+1)-(x^{2}-1)\times 2x}{(x^{2}+1)^{2}} &&\\
k'(x) &= \dfrac{2x^{3}+2x-2x^{3}+2x}{(x^{2}+1)^{2}} &&\justif{$-(x^{2}-1)\times 2x=-2x^{3}+2x$}\\
k'(x) &= \res{\dfrac{4x}{(x^{2}+1)^{2}}} &&
\end{calculs}
\rem{Le \og$-$\fg{} devant $uv'$ porte sur tout le produit : on garde $uv'$ entre
parenthèses jusqu'au développement. Le dénominateur $v^{2}$ ne se développe pas.}
\end{exo}

\begin{exo}[Entraînement : inverses et quotients]
\cas{a)}
\begin{calculs}
v(x) &= 3-x &&\justif{$v'(x)=-1$}\\
f'(x) &= -\dfrac{-1}{(3-x)^{2}} &&\justif{$-\frac{v'}{v^{2}}$}\\
f'(x) &= \res{\dfrac{1}{(3-x)^{2}}} &&\justif{deux signes \og$-$\fg}
\end{calculs}
\cas{b)}
\begin{calculs}
g(x) &= -2\times\dfrac{1}{x^{2}+x+1} &&\justif{le coefficient $-2$ se garde}\\
v(x) &= x^{2}+x+1 &&\justif{$v'(x)=2x+1$}\\
g'(x) &= -2\times\left(-\dfrac{2x+1}{(x^{2}+x+1)^{2}}\right) &&\\
g'(x) &= \res{\dfrac{4x+2}{(x^{2}+x+1)^{2}}} &&
\end{calculs}
\cas{c)}
\begin{calculs}
u(x) &= 5x+2 &&\justif{$u'(x)=5$}\\
v(x) &= x-1 &&\justif{$v'(x)=1$}\\
h'(x) &= \dfrac{5(x-1)-(5x+2)\times 1}{(x-1)^{2}} &&\justif{$\frac{u'v-uv'}{v^{2}}$}\\
h'(x) &= \dfrac{5x-5-5x-2}{(x-1)^{2}} &&\\
h'(x) &= \res{\dfrac{-7}{(x-1)^{2}}} &&
\end{calculs}
\cas{d)}
\begin{calculs}
u(x) &= 1-2x &&\justif{$u'(x)=-2$}\\
v(x) &= x+3 &&\justif{$v'(x)=1$}\\
k'(x) &= \dfrac{-2(x+3)-(1-2x)\times 1}{(x+3)^{2}} &&\\
k'(x) &= \dfrac{-2x-6-1+2x}{(x+3)^{2}} &&\justif{$-(1-2x)=-1+2x$}\\
k'(x) &= \res{\dfrac{-7}{(x+3)^{2}}} &&
\end{calculs}
\cas{e)}
\begin{calculs}
u(x) &= x^{2}+1 &&\justif{$u'(x)=2x$}\\
v(x) &= x &&\justif{$v'(x)=1$}\\
m'(x) &= \dfrac{2x\times x-(x^{2}+1)\times 1}{x^{2}} &&\\
m'(x) &= \dfrac{2x^{2}-x^{2}-1}{x^{2}} &&\\
m'(x) &= \res{\dfrac{x^{2}-1}{x^{2}}} &&
\end{calculs}
\rem{Plus court en e) : $m(x)=x+\dfrac{1}{x}$, de dérivée $1-\dfrac{1}{x^{2}}$, la même
expression.}
\cas{f)}
\begin{calculs}
u(x) &= 3x &&\justif{$u'(x)=3$}\\
v(x) &= x^{2}+1 &&\justif{$v'(x)=2x$}\\
p'(x) &= \dfrac{3(x^{2}+1)-3x\times 2x}{(x^{2}+1)^{2}} &&\\
p'(x) &= \dfrac{3x^{2}+3-6x^{2}}{(x^{2}+1)^{2}} &&\\
p'(x) &= \res{\dfrac{3-3x^{2}}{(x^{2}+1)^{2}}} &&
\end{calculs}
\cas{g)}
\begin{calculs}
u(x) &= x^{2}-4x &&\justif{$u'(x)=2x-4$}\\
v(x) &= x+2 &&\justif{$v'(x)=1$}\\
q'(x) &= \dfrac{(2x-4)(x+2)-(x^{2}-4x)\times 1}{(x+2)^{2}} &&\\
q'(x) &= \dfrac{2x^{2}-8-x^{2}+4x}{(x+2)^{2}} &&\justif{$(2x-4)(x+2)=2x^{2}+4x-4x-8$}\\
q'(x) &= \res{\dfrac{x^{2}+4x-8}{(x+2)^{2}}} &&
\end{calculs}
\cas{h)}
\begin{calculs}
u(x) &= \sqrt{x} &&\justif{$u'(x)=\frac{1}{2\sqrt{x}}$}\\
v(x) &= x+1 &&\justif{$v'(x)=1$}\\
r'(x) &= \dfrac{\frac{1}{2\sqrt{x}}\times(x+1)-\sqrt{x}\times 1}{(x+1)^{2}} &&\justif{$\frac{u'v-uv'}{v^{2}}$}\\
r'(x) &= \dfrac{(x+1)-2\sqrt{x}\times\sqrt{x}}{2\sqrt{x}\,(x+1)^{2}} &&\justif{on multiplie en haut et en bas par $2\sqrt{x}$}\\
r'(x) &= \dfrac{x+1-2x}{2\sqrt{x}\,(x+1)^{2}} &&\justif{$\sqrt{x}\times\sqrt{x}=x$}\\
r'(x) &= \res{\dfrac{1-x}{2\sqrt{x}\,(x+1)^{2}}} &&
\end{calculs}
\explique{$r$ est définie pour $x\ge 0$, mais dérivable seulement pour $x>0$ : la racine
carrée n'est pas dérivable en $0$.}{propriété 5}
\end{exo}

\begin{exo}[Choisir la bonne formule]
\explique{Pour reconnaître la forme, on cherche la \emph{dernière} opération du calcul de
$f(x)$.}{en c), on calcule $x^{2}+3$, puis on divise $4$ par le résultat}
\cas{a)}
\explique{$f(x)=3x^{4}-\dfrac{2}{x}$ : une \textbf{somme} (différence de deux
termes).}{on dérive chaque terme}
\begin{calculs}
f'(x) &= 3\times 4x^{3}-2\times\left(-\dfrac{1}{x^{2}}\right) &&\\
f'(x) &= \res{12x^{3}+\dfrac{2}{x^{2}}} &&
\end{calculs}
\cas{b)}
\explique{$g(x)=(x-1)(x^{2}+2)$ : un \textbf{produit}.}{$u'v+uv'$}
\begin{calculs}
u(x) &= x-1 &&\justif{$u'(x)=1$}\\
v(x) &= x^{2}+2 &&\justif{$v'(x)=2x$}\\
g'(x) &= 1\times(x^{2}+2)+(x-1)\times 2x &&\\
g'(x) &= x^{2}+2+2x^{2}-2x &&\\
g'(x) &= \res{3x^{2}-2x+2} &&
\end{calculs}
\cas{c)}
\explique{$h(x)=4\times\dfrac{1}{x^{2}+3}$ : un \textbf{inverse}, multiplié par
$4$.}{le numérateur est un nombre : pas besoin du quotient}
\begin{calculs}
v(x) &= x^{2}+3 &&\justif{$v'(x)=2x$}\\
h'(x) &= 4\times\left(-\dfrac{2x}{(x^{2}+3)^{2}}\right) &&\justif{$4\times\left(-\frac{v'}{v^{2}}\right)$}\\
h'(x) &= \res{-\dfrac{8x}{(x^{2}+3)^{2}}} &&
\end{calculs}
\cas{d)}
\explique{$k(x)=\dfrac{x-4}{2x+1}$ : un \textbf{quotient}.}{$\frac{u'v-uv'}{v^{2}}$}
\begin{calculs}
u(x) &= x-4 &&\justif{$u'(x)=1$}\\
v(x) &= 2x+1 &&\justif{$v'(x)=2$}\\
k'(x) &= \dfrac{1\times(2x+1)-(x-4)\times 2}{(2x+1)^{2}} &&\\
k'(x) &= \dfrac{2x+1-2x+8}{(2x+1)^{2}} &&\justif{$-(x-4)\times 2=-2x+8$}\\
k'(x) &= \res{\dfrac{9}{(2x+1)^{2}}} &&
\end{calculs}
\cas{e)}
\explique{$m(x)=x^{2}\sqrt{x}$ : un \textbf{produit}.}{$u(x)=x^{2}$ et $v(x)=\sqrt{x}$}
\begin{calculs}
m'(x) &= 2x\sqrt{x}+x^{2}\times\dfrac{1}{2\sqrt{x}} &&\justif{$u'v+uv'$}\\
m'(x) &= 2x\sqrt{x}+\dfrac{x\sqrt{x}}{2} &&\justif{$\frac{x^{2}}{\sqrt{x}}=\frac{x\times\sqrt{x}\times\sqrt{x}}{\sqrt{x}}=x\sqrt{x}$}\\
m'(x) &= \res{\dfrac{5}{2}x\sqrt{x}} &&\justif{$2+\frac{1}{2}=\frac{5}{2}$}
\end{calculs}
\cas{f)}
\explique{$p(x)=\sqrt{x}+\sin x$ : une \textbf{somme}.}{}
\begin{calculs}
p'(x) &= \res{\dfrac{1}{2\sqrt{x}}+\cos x} &&
\end{calculs}
\end{exo}

\begin{exo}[Tangentes à une hyperbole]
\cas{1.}
\begin{calculs}
u(x) &= 2x-4 &&\justif{$u'(x)=2$}\\
v(x) &= x+1 &&\justif{$v'(x)=1$}\\
f'(x) &= \dfrac{2(x+1)-(2x-4)\times 1}{(x+1)^{2}} &&\justif{$\frac{u'v-uv'}{v^{2}}$}\\
f'(x) &= \dfrac{2x+2-2x+4}{(x+1)^{2}} &&\\
f'(x) &= \res{\dfrac{6}{(x+1)^{2}}} &&
\end{calculs}
\cas{2.}
\explique{$A$ est sur l'axe des abscisses : son ordonnée est nulle, on résout
$f(x)=0$.}{une fraction est nulle quand son numérateur est nul}
\begin{calculs}
2x-4 &= 0 &&\\
x &= 2 &&\justif{$2\ne -1$ : ce n'est pas la valeur interdite}
\end{calculs}
\explique{Donc $\res{A(2\,;0)}$.}{}
\begin{calculs}
f'(2) &= \dfrac{6}{(2+1)^{2}} = \dfrac{6}{9} &&\\
f'(2) &= \dfrac{2}{3} &&\justif{on simplifie par $3$}\\
y &= f'(2)(x-2)+f(2) &&\justif{propriété 1}\\
y &= \dfrac{2}{3}(x-2)+0 &&\\
y &= \res{\dfrac{2}{3}x-\dfrac{4}{3}} &&
\end{calculs}
\cas{3.}
\explique{$B$ est sur l'axe des ordonnées : son abscisse est nulle, on calcule
$f(0)$.}{}
\begin{calculs}
f(0) &= \dfrac{0-4}{0+1} = -4 &&\justif{donc $\res{B(0\,;-4)}$}\\
f'(0) &= \dfrac{6}{(0+1)^{2}} = 6 &&\\
y &= f'(0)(x-0)+f(0) &&\\
y &= \res{6x-4} &&
\end{calculs}
\rem{Ne pas confondre : sur l'axe des abscisses, $y=0$ ; sur l'axe des ordonnées, $x=0$.}
\end{exo}

\begin{exo}[Deux fonctions, une même dérivée]
\cas{1.}
\begin{calculs}
u(x) &= 2x+1 &&\justif{$u'(x)=2$}\\
v(x) &= x-1 &&\justif{$v'(x)=1$}\\
f'(x) &= \dfrac{2(x-1)-(2x+1)\times 1}{(x-1)^{2}} &&\justif{$\frac{u'v-uv'}{v^{2}}$}\\
f'(x) &= \dfrac{2x-2-2x-1}{(x-1)^{2}} &&\\
f'(x) &= \res{\dfrac{-3}{(x-1)^{2}}} &&
\end{calculs}
\begin{calculs}
g(x) &= 3\times\dfrac{1}{x-1} &&\justif{un inverse, multiplié par $3$}\\
g'(x) &= 3\times\left(-\dfrac{1}{(x-1)^{2}}\right) &&\justif{$v(x)=x-1$, $v'(x)=1$}\\
g'(x) &= \res{\dfrac{-3}{(x-1)^{2}}} &&
\end{calculs}
\explique{On remarque que $f$ et $g$ ont \res{la même dérivée}.}{}
\cas{2.}
\begin{calculs}
f(x)-g(x) &= \dfrac{2x+1}{x-1}-\dfrac{3}{x-1} &&\justif{même dénominateur}\\
f(x)-g(x) &= \dfrac{2x+1-3}{x-1} &&\\
f(x)-g(x) &= \dfrac{2\barre{(x-1)}}{\barre{(x-1)}} &&\justif{$2x-2=2(x-1)$}\\
f(x)-g(x) &= \res{2} &&
\end{calculs}
\explique{$f(x)=g(x)+2$ : la courbe de $f$ est celle de $g$ montée de $2$. Ses tangentes
sont montées de $2$ elles aussi, sans changer de pente : d'où la même
dérivée.}{$(g+2)'=g'+0$ : la dérivée d'une constante est nulle}
\end{exo}

\partie{Dérivée des fonctions composées}

\begin{exo}[Dériver une puissance, une racine]
\explique{Dans chaque cas, on repère la fonction $u$ \og à l'intérieur\fg, on calcule $u'$,
puis on applique la ligne du tableau de la propriété 8.}{$(u^{2})'=2u'u$ ;
$(u^{n})'=nu'u^{n-1}$ ; $(\sqrt{u})'=\frac{u'}{2\sqrt{u}}$}
\cas{a)}
\begin{calculs}
u(x) &= 3x+1 &&\justif{$u'(x)=3$}\\
f'(x) &= 2\times 3\times(3x+1) &&\justif{$2u'u$}\\
f'(x) &= \res{18x+6} &&
\end{calculs}
\explique{Vérification : $f(x)=9x^{2}+6x+1$ en développant, de dérivée $18x+6$.}{partie
V}
\cas{b)}
\begin{calculs}
u(x) &= 1-2x &&\justif{$u'(x)=-2$}\\
g'(x) &= 2\times(-2)\times(1-2x) &&\justif{$2u'u$}\\
g'(x) &= -4+8x &&\\
g'(x) &= \res{8x-4} &&
\end{calculs}
\cas{c)}
\begin{calculs}
u(x) &= x^{2}+3 &&\justif{$u'(x)=2x$}\\
h'(x) &= 3\times 2x\times(x^{2}+3)^{2} &&\justif{$nu'u^{n-1}$, avec $n=3$}\\
h'(x) &= \res{6x(x^{2}+3)^{2}} &&
\end{calculs}
\cas{d)}
\begin{calculs}
u(x) &= 4x-1 &&\justif{$u'(x)=4$}\\
k'(x) &= 5\times 4\times(4x-1)^{4} &&\justif{$nu'u^{n-1}$, avec $n=5$}\\
k'(x) &= \res{20(4x-1)^{4}} &&
\end{calculs}
\cas{e)}
\begin{calculs}
u(x) &= 2x+6 &&\justif{$u'(x)=2$ ; $u(x)>0$ pour $x>-3$}\\
m'(x) &= \dfrac{2}{2\sqrt{2x+6}} &&\justif{$\frac{u'}{2\sqrt{u}}$}\\
m'(x) &= \res{\dfrac{1}{\sqrt{2x+6}}} &&\justif{on simplifie par $2$}
\end{calculs}
\cas{f)}
\begin{calculs}
u(x) &= x^{2}+4 &&\justif{$u'(x)=2x$ ; $u(x)\ge 4>0$ sur $\R$}\\
p'(x) &= \dfrac{2x}{2\sqrt{x^{2}+4}} &&\justif{$\frac{u'}{2\sqrt{u}}$}\\
p'(x) &= \res{\dfrac{x}{\sqrt{x^{2}+4}}} &&
\end{calculs}
\rem{Erreur fréquente : oublier le facteur $u'$ ($2(3x+1)$ au lieu de $6(3x+1)$). En c) et
d), on ne développe pas : la forme $6x(x^{2}+3)^{2}$ est plus simple.}
\end{exo}

\begin{exo}[Cosinus, sinus et formule générale]
\cas{1.}
\explique{On garde $u(x)$ à l'intérieur du sinus ou du cosinus, et on multiplie par
$u'(x)$.}{$(\sin u)'=u'\cos u$ ; $(\cos u)'=-u'\sin u$}
\cas{a)}
\begin{calculs}
u(x) &= 3x &&\justif{$u'(x)=3$}\\
f'(x) &= \res{3\cos(3x)} &&\justif{$u'\cos u$}
\end{calculs}
\cas{b)}
\begin{calculs}
u(x) &= 5x+1 &&\justif{$u'(x)=5$}\\
g'(x) &= \res{-5\sin(5x+1)} &&\justif{$-u'\sin u$}
\end{calculs}
\cas{c)}
\begin{calculs}
u(x) &= x^{2} &&\justif{$u'(x)=2x$}\\
h'(x) &= \res{2x\cos(x^{2})} &&\justif{$u'\cos u$}
\end{calculs}
\cas{d)}
\begin{calculs}
u(x) &= 1-x &&\justif{$u'(x)=-1$}\\
k'(x) &= -(-1)\sin(1-x) &&\justif{$-u'\sin u$}\\
k'(x) &= \res{\sin(1-x)} &&\justif{deux signes \og$-$\fg}
\end{calculs}
\rem{L'intérieur reste tel quel : $\cos(5x+1)$ donne $-5\sin(5x+1)$, et non $-\sin(5)$ ou
$-\sin(5x+1)$ sans le $5$.}
\cas{2.}
\explique{$F(x)=G(2x)$ : la fonction \og intérieure\fg{} est $u(x)=2x$, et $G$ est la
fonction \og extérieure\fg.}{$F(x)=G\big(u(x)\big)$}
\begin{calculs}
u(x) &= 2x &&\justif{$u'(x)=2$}\\
G'\big(u(x)\big) &= \dfrac{1}{1+(2x)^{2}} &&\justif{on remplace $x$ par $2x$ dans $G'(x)$}\\
F'(x) &= u'(x)\times G'\big(u(x)\big) &&\justif{on recopie la formule}\\
F'(x) &= 2\times\dfrac{1}{1+(2x)^{2}} &&\\
F'(x) &= \res{\dfrac{2}{1+4x^{2}}} &&\justif{$(2x)^{2}=4x^{2}$}
\end{calculs}
\rem{Erreur fréquente : écrire $(2x)^{2}=2x^{2}$ ; le carré porte aussi sur le $2$.}
\end{exo}

\begin{exo}[Composée et tangente]
\cas{1.}
\begin{calculs}
u(x) &= x^{2}+9 &&\justif{$u'(x)=2x$ ; $u(x)\ge 9>0$ sur $\R$}\\
f'(x) &= \dfrac{2x}{2\sqrt{x^{2}+9}} &&\justif{$\frac{u'}{2\sqrt{u}}$}\\
f'(x) &= \res{\dfrac{x}{\sqrt{x^{2}+9}}} &&\justif{on simplifie par $2$}
\end{calculs}
\cas{2.}
\begin{calculs}
f(4) &= \sqrt{16+9} = \sqrt{25} &&\\
f(4) &= \res{5} &&\\
f'(4) &= \dfrac{4}{\sqrt{25}} = \dfrac{4}{5} &&\\
f'(4) &= \res{0{,}8} &&
\end{calculs}
\cas{3.}
\begin{calculs}
y &= f'(4)(x-4)+f(4) &&\justif{propriété 1}\\
y &= 0{,}8(x-4)+5 &&\\
y &= 0{,}8x-3{,}2+5 &&\\
y &= \res{0{,}8x+1{,}8} &&
\end{calculs}
\cas{4.}
\explique{L'approximation linéaire, c'est l'équation de la tangente avec $x=4{,}1$ :
$a=4$ et $h=0{,}1$.}{propriété 2}
\begin{calculs}
f(4+0{,}1) &\approx f(4)+f'(4)\times 0{,}1 &&\justif{on recopie la formule}\\
f(4{,}1) &\approx 5+0{,}8\times 0{,}1 &&\\
f(4{,}1) &\approx \res{5{,}08} &&
\end{calculs}
\explique{La calculatrice donne $f(4{,}1)=\sqrt{25{,}81}\approx 5{,}0804$ : la valeur
approchée est juste au centième.}{écart d'environ $0{,}0004$}
\end{exo}

\partie{Étude des variations d'une fonction}

\begin{exo}[Du signe de $f'$ aux variations de $f$]
\cas{1.}
\explique{Là où $f'(x)\le 0$, $f$ est décroissante ; là où $f'(x)\ge 0$, $f$ est
croissante.}{théorème 1}
\explique{On place les images de l'énoncé aux extrémités des flèches : $f$ décroît de $2$ à
$-3$, croît de $-3$ à $4$, puis décroît de $4$ à $1$.}{une flèche va d'une valeur à la
suivante}
\begin{center}
\begin{tikzpicture}[scale=0.9]
\tkzTabInit[lgt=2,espcl=2.2]{$x$/0.8,$f'(x)$/0.8,$f(x)$/1.6}{$-4$,$-1$,$3$,$5$}
\tkzTabLine{,-,z,+,z,-,}
\tkzTabVar{+/$2$,-/$-3$,+/$4$,-/$1$}
\end{tikzpicture}
\end{center}
\cas{2.}
\explique{$0$ et $2$ sont tous deux dans $\intff{-1}{3}$, où $f$ est croissante. Comme
$0<2$, on a $\res{f(0)<f(2)}$.}{croissante : les images sont rangées dans le même ordre
que les nombres}
\rem{On ne connaît ni $f(0)$ ni $f(2)$ : seul le sens de variation permet de les
comparer, et il faut que les deux nombres soient dans le même intervalle de monotonie.
L'inégalité est stricte, car $f'(x)>0$ entre $-1$ et $3$.}
\end{exo}

\begin{exo}[Variations de polynômes]
\explique{Pour chaque fonction : on dérive, on cherche où la dérivée s'annule, on étudie
son signe, puis on calcule les images des valeurs trouvées.}{méthode 11}
\cas{a)}
\begin{calculs}
f'(x) &= 2x-4 &&\justif{on dérive}\\
2x-4 &= 0 &&\justif{on cherche où $f'$ s'annule}\\
x &= 2 &&\\
f(2) &= 4-8+7 = 3 &&\justif{image de $2$}
\end{calculs}
\explique{$2x-4$ est négatif avant $2$, positif après.}{affine, coefficient $2>0$ : $-$
puis $+$}
\begin{center}
\begin{tikzpicture}[scale=0.9]
\tkzTabInit[lgt=2,espcl=2.4]{$x$/0.8,$f'(x)$/0.8,$f(x)$/1.6}{$-\infty$,$2$,$+\infty$}
\tkzTabLine{,-,z,+,}
\tkzTabVar{+/{},-/$3$,+/{}}
\end{tikzpicture}
\end{center}
\explique{On retrouve le sommet de la parabole, d'abscisse
$-\dfrac{b}{2a}=\dfrac{4}{2}=2$.}{contrôle avec le second degré}
\cas{b)}
\begin{calculs}
g'(x) &= -6x+12 &&\justif{on dérive}\\
-6x+12 &= 0 &&\\
x &= 2 &&\\
g(2) &= -12+24-5 = 7 &&
\end{calculs}
\explique{$-6x+12$ est positif avant $2$, négatif après.}{coefficient $-6<0$ : $+$ puis
$-$}
\begin{center}
\begin{tikzpicture}[scale=0.9]
\tkzTabInit[lgt=2,espcl=2.4]{$x$/0.8,$g'(x)$/0.8,$g(x)$/1.6}{$-\infty$,$2$,$+\infty$}
\tkzTabLine{,+,z,-,}
\tkzTabVar{-/{},+/$7$,-/{}}
\end{tikzpicture}
\end{center}
\cas{c)}
\begin{calculs}
h'(x) &= 6x^{2}-6x-12 &&\justif{on dérive}
\end{calculs}
\explique{Racines du trinôme $6x^{2}-6x-12$ :}{on cherche où $h'$ s'annule}
\begin{calculs}
\Delta &= (-6)^{2}-4\times 6\times(-12) &&\\
\Delta &= 36+288 = 324 &&\justif{$\sqrt{324}=18$}\\
x_{1} &= \dfrac{6-18}{12} = -1 &&\\
x_{2} &= \dfrac{6+18}{12} = 2 &&
\end{calculs}
\explique{Images des racines :}{les nombres de la ligne des $x$}
\begin{calculs}
h(-1) &= -2-3+12+1 = 8 &&\justif{$2\times(-1)^{3}=-2$}\\
h(2) &= 16-12-24+1 = -19 &&
\end{calculs}
\explique{$h'(x)$ est du signe de $a=6$, sauf entre les racines.}{signe du trinôme}
\begin{center}
\begin{tikzpicture}[scale=0.9]
\tkzTabInit[lgt=2,espcl=2.2]{$x$/0.8,$h'(x)$/0.8,$h(x)$/1.6}{$-\infty$,$-1$,$2$,$+\infty$}
\tkzTabLine{,+,z,-,z,+,}
\tkzTabVar{-/{},+/$8$,-/$-19$,+/{}}
\end{tikzpicture}
\end{center}
\rem{On pouvait aussi factoriser : $h'(x)=6(x^{2}-x-2)=6(x+1)(x-2)$, avec la racine
évidente $-1$.}
\cas{d)}
\begin{calculs}
k'(x) &= 3x^{2}+3 &&\justif{on dérive}
\end{calculs}
\explique{$x^{2}\ge 0$, donc $3x^{2}+3\ge 3>0$ : $k'(x)$ est toujours strictement
positif.}{inutile de calculer $\Delta$}
\begin{center}
\begin{tikzpicture}[scale=0.9]
\tkzTabInit[lgt=2,espcl=3]{$x$/0.8,$k'(x)$/0.8,$k(x)$/1.6}{$-\infty$,$+\infty$}
\tkzTabLine{,+,}
\tkzTabVar{-/{},+/{}}
\end{tikzpicture}
\end{center}
\explique{$k$ est \res{strictement croissante sur $\R$}.}{théorème 1}
\end{exo}

\begin{exo}[Variations d'un quotient]
\cas{1.}
\begin{calculs}
u(x) &= x^{2}+3 &&\justif{$u'(x)=2x$}\\
v(x) &= x+1 &&\justif{$v'(x)=1$}\\
f'(x) &= \dfrac{2x(x+1)-(x^{2}+3)\times 1}{(x+1)^{2}} &&\justif{$\frac{u'v-uv'}{v^{2}}$}\\
f'(x) &= \dfrac{2x^{2}+2x-x^{2}-3}{(x+1)^{2}} &&\\
f'(x) &= \res{\dfrac{x^{2}+2x-3}{(x+1)^{2}}} &&
\end{calculs}
\cas{2.}
\explique{Sur $\intoo{-1}{+\infty}$, $(x+1)^{2}>0$ : $f'(x)$ a le signe du numérateur
$x^{2}+2x-3$.}{un carré non nul est strictement positif}
\begin{calculs}
\Delta &= 2^{2}-4\times 1\times(-3) = 16 &&\justif{$\sqrt{16}=4$}\\
x_{1} &= \dfrac{-2-4}{2} = -3 &&\\
x_{2} &= \dfrac{-2+4}{2} = 1 &&
\end{calculs}
\explique{Le trinôme ($a=1>0$) est négatif entre $-3$ et $1$, positif à l'extérieur. Sur
$\intoo{-1}{+\infty}$, seule la racine $1$ compte : $\res{f'(x)<0}$ sur
$\intoo{-1}{1}$ et $\res{f'(x)>0}$ sur $\intoo{1}{+\infty}$.}{$-3$ n'est pas dans
l'intervalle d'étude}
\cas{3.}
\begin{calculs}
f(1) &= \dfrac{1+3}{1+1} = 2 &&\justif{image de la racine de $f'$}
\end{calculs}
\begin{center}
\begin{tikzpicture}[scale=0.9]
\tkzTabInit[lgt=2,espcl=2.4]{$x$/0.8,$f'(x)$/0.8,$f(x)$/1.6}{$-1$,$1$,$+\infty$}
\tkzTabLine{d,-,z,+,}
\tkzTabVar{D+/{},-/$2$,+/{}}
\end{tikzpicture}
\end{center}
\explique{La double barre en $-1$ rappelle que $f$ n'y est pas définie.}{valeur interdite}
\end{exo}

\begin{exo}[Lire le signe de la dérivée]
\begin{center}
\begin{tikzpicture}[x=0.8cm,y=0.8cm]
  \repere{-5}{-2}{5}{4}
  \gradx{-4,-3,-2,-1,1,2,3,4} \grady{-1,1,2,3}
  \draw[coulCourbe,line width=1.1pt]
    (-4,2) .. controls (-3.333,0.667) and (-2.667,-1) .. (-2,-1)
           .. controls (-1,-1) and (0,3) .. (1,3)
           .. controls (1.667,3) and (2.333,1) .. (3,1)
           .. controls (3.333,1) and (3.667,1.333) .. (4,2);
  \draw[coulExo,line width=0.9pt] (-2.8,-1) -- (-1.2,-1) (0.2,3) -- (1.8,3)
    (2.3,1) -- (3.7,1);
  \draw[coulRep,dashed,thin] (-2,0) -- (-2,-1) (1,0) -- (1,3) (3,0) -- (3,1);
  \path[coulCourbe] (-2,-1) node[croix]{} (1,3) node[croix]{} (3,1) node[croix]{};
  \node[coulRep,font=\small] at (-3,-1.6) {$-$};
  \node[coulRep,font=\small] at (-0.5,-1.6) {$+$};
  \node[coulRep,font=\small] at (2,-1.6) {$-$};
  \node[coulRep,font=\small] at (3.5,-1.6) {$+$};
  \node[coulCourbe,font=\footnotesize] at (-4.4,2.8) {$\Cg$};
\end{tikzpicture}
\end{center}
\cas{1.}
\explique{Là où la courbe descend, ses tangentes descendent : $g'(x)\le 0$. Là où elle
monte, $g'(x)\ge 0$.}{le signe de $g'(a)$ est celui de la pente de la tangente en $a$}
\explique{Aux tangentes horizontales, en $-2$, en $1$ et en $3$, on a $g'(x)=0$.}{signes
reportés sous l'axe de la figure}
\begin{center}
\begin{tikzpicture}[scale=0.9]
\tkzTabInit[lgt=2,espcl=1.9]{$x$/0.8,$g'(x)$/0.8}{$-4$,$-2$,$1$,$3$,$4$}
\tkzTabLine{,-,z,+,z,-,z,+,}
\end{tikzpicture}
\end{center}
\cas{2.}
\explique{Les images se lisent sur la courbe : $g(-4)=2$, $g(-2)=-1$, $g(1)=3$, $g(3)=1$
et $g(4)=2$.}{ordonnées des points de la figure}
\begin{center}
\begin{tikzpicture}[scale=0.9]
\tkzTabInit[lgt=2,espcl=1.9]{$x$/0.8,$g'(x)$/0.8,$g(x)$/1.6}{$-4$,$-2$,$1$,$3$,$4$}
\tkzTabLine{,-,z,+,z,-,z,+,}
\tkzTabVar{+/$2$,-/$-1$,+/$3$,-/$1$,+/$2$}
\end{tikzpicture}
\end{center}
\rem{Ne pas confondre le signe de $g(x)$ (la courbe est au-dessus ou au-dessous de l'axe)
et le signe de $g'(x)$ (la courbe monte ou descend). Sur $\intff{-4}{-2}$, $g(x)$
change de signe alors que $g'(x)$ reste négatif.}
\end{exo}

\partie{Extremums d'une fonction}

\begin{exo}[Lire les extremums]
\cas{1.}
\explique{On compare les \og sommets\fg{} du tableau, bornes comprises : $1$ ; $5$ ; $-3$ et
$0$.}{le maximum est la plus grande valeur prise par $f$}
\explique{Le maximum de $f$ sur $\intff{-2}{6}$ est $\res{5}$, atteint en $1$ ; le minimum
est $\res{-3}$, atteint en $4$.}{$f(-2)=1$ et $f(6)=0$ restent entre les deux}
\cas{2.}
\explique{$f$ est croissante sur $\intff{-2}{1}$ : $\res{f'(x)\ge 0}$ ; décroissante sur
$\intff{1}{4}$ : $\res{f'(x)\le 0}$ ; croissante sur $\intff{4}{6}$ :
$\res{f'(x)\ge 0}$.}{courbe qui monte : tangentes qui montent, coefficient directeur
positif}
\cas{3.}
\explique{En $1$, $f$ passe de croissante à décroissante : maximum local, la tangente y est
horizontale, donc $f'(1)=\res{0}$. En $4$ (minimum local), de même :
$f'(4)=\res{0}$.}{figures du théorème 2}
\end{exo}

\begin{exo}[Extremums d'un polynôme]
\cas{1.}
\begin{calculs}
f'(x) &= 3x^{2}-12x+9 &&\justif{on dérive}\\
f'(x) &= 3(x^{2}-4x+3) &&\justif{on factorise par $3$}
\end{calculs}
\explique{Racines de $x^{2}-4x+3$ :}{on cherche où $f'$ s'annule}
\begin{calculs}
\Delta &= (-4)^{2}-4\times 1\times 3 = 4 &&\justif{$\sqrt{4}=2$}\\
x_{1} &= \dfrac{4-2}{2} = 1 &&\\
x_{2} &= \dfrac{4+2}{2} = 3 &&
\end{calculs}
\explique{$f'(x)$ est du signe de $3>0$ à l'extérieur de $1$ et $3$, négatif
entre.}{signe du trinôme}
\explique{Images des racines et des bornes :}{les nombres de la ligne des $x$}
\begin{calculs}
f(0) &= 1 &&\\
f(1) &= 1-6+9+1 = 5 &&\\
f(3) &= 27-54+27+1 = 1 &&\\
f(5) &= 125-150+45+1 = 21 &&
\end{calculs}
\begin{center}
\begin{tikzpicture}[scale=0.9]
\tkzTabInit[lgt=2,espcl=2.2]{$x$/0.8,$f'(x)$/0.8,$f(x)$/1.6}{$0$,$1$,$3$,$5$}
\tkzTabLine{,+,z,-,z,+,}
\tkzTabVar{-/$1$,+/$5$,-/$1$,+/$21$}
\end{tikzpicture}
\end{center}
\cas{2.}
\explique{Le maximum de $f$ sur $\intff{0}{5}$ est $\res{21}$, atteint en $5$ ; le minimum
est $\res{1}$, atteint en $0$ et en $3$.}{on compare les sommets, bornes comprises : $1$ ;
$5$ ; $1$ et $21$}
\cas{3.}
\explique{Oui : $f'$ s'annule en $1$ en passant de $+$ à $-$, donc $f$ a un \res{maximum
local égal à $5$} en $1$.}{théorème 2}
\rem{Ce maximum local n'est pas le maximum sur $\intff{0}{5}$, car $f(5)=21>5$.
\og Local\fg{} veut dire : la plus grande valeur seulement autour de $1$.}
\end{exo}

\begin{exo}[Vrai ou faux]
\cas{1.}
\explique{\res{Faux.} Contre-exemple : $f(x)=x^{3}-6x^{2}+12x$.}{un seul contre-exemple
suffit pour montrer qu'une affirmation est fausse}
\begin{calculs}
f'(x) &= 3x^{2}-12x+12 &&\justif{on dérive}\\
f'(x) &= 3(x^{2}-4x+4) &&\\
f'(x) &= 3(x-2)^{2} &&\justif{$(a-b)^{2}=a^{2}-2ab+b^{2}$}
\end{calculs}
\explique{$f'(2)=0$, mais $3(x-2)^{2}\ge 0$ : $f'$ ne change pas de signe, $f$ est
croissante sur $\R$, sans extremum en $2$.}{il faut que $f'$ \emph{change de signe}
(théorème 2)}
\rem{C'est le contre-exemple du cours ($x^{3}$ en $0$), décalé pour que la dérivée
s'annule en $2$.}
\cas{2.}
\begin{calculs}
g'(x) &= 3x^{2}+1 &&\justif{on dérive}
\end{calculs}
\explique{$3x^{2}\ge 0$, donc $g'(x)\ge 1>0$ : $g$ est strictement croissante sur $\R$.
Elle n'a aucune tangente horizontale, donc aucun extremum.}{$g'$ ne s'annule jamais}
\explique{L'affirmation est \res{vraie}.}{}
\cas{3.}
\begin{calculs}
h'(x) &= 1-\dfrac{1}{x^{2}} &&\justif{$x$ donne $1$ ; $\frac{1}{x}$ donne $-\frac{1}{x^{2}}$}\\
h'(x) &= \dfrac{1\fois{x^{2}}}{1\fois{x^{2}}}-\dfrac{1}{x^{2}} &&\justif{même dénominateur}\\
h'(x) &= \dfrac{x^{2}-1}{x^{2}} &&\\
h'(x) &= \dfrac{(x-1)(x+1)}{x^{2}} &&\justif{$a^{2}-b^{2}=(a-b)(a+b)$}
\end{calculs}
\explique{Sur $\intoo{0}{+\infty}$, $x^{2}>0$ et $x+1>0$ : $h'(x)$ a le signe de
$x-1$.}{on ne garde que le facteur dont le signe change}
\begin{calculs}
h(1) &= 1+\dfrac{1}{1} = 2 &&
\end{calculs}
\begin{center}
\begin{tikzpicture}[scale=0.9]
\tkzTabInit[lgt=2,espcl=2.4]{$x$/0.8,$h'(x)$/0.8,$h(x)$/1.6}{$0$,$1$,$+\infty$}
\tkzTabLine{d,-,z,+,}
\tkzTabVar{D+/{},-/$2$,+/{}}
\end{tikzpicture}
\end{center}
\explique{$h$ décroît puis croît : son minimum sur $\intoo{0}{+\infty}$ est $h(1)=2$.
L'affirmation est \res{vraie}.}{}
\end{exo}

\partie{Étude complète d'une fonction}

\begin{exo}[Étude d'un polynôme]
\cas{1.}
\explique{$f$ est un polynôme : pas de valeur interdite. Elle est \res{définie et dérivable
sur $\R$}, donc sur $\intff{-1}{3}$.}{propriété 4 : un polynôme se dérive partout}
\cas{2.}
\begin{calculs}
f'(x) &= -3x^{2}+3\times 2x &&\justif{la constante $-1$ donne $0$}\\
f'(x) &= -3x^{2}+6x &&\\
f'(x) &= \res{-3x(x-2)} &&\justif{on factorise par $-3x$}
\end{calculs}
\explique{Vérification : $-3x\times x=-3x^{2}$ et $-3x\times(-2)=6x$.}{on redéveloppe}
\cas{3.}
\explique{On étudie le signe de chaque facteur sur $\intff{-1}{3}$ : $-3x$ est positif
avant $0$, négatif après ; $x-2$ est négatif avant $2$, positif après.}{tableau de signes
d'un produit}
\begin{center}
\begin{tikzpicture}[scale=0.9]
\tkzTabInit[lgt=2,espcl=2.2]{$x$/0.8,$-3x$/0.8,$x-2$/0.8,$f'(x)$/0.8}{$-1$,$0$,$2$,$3$}
\tkzTabLine{,+,z,-,t,-,}
\tkzTabLine{,-,t,-,z,+,}
\tkzTabLine{,-,z,+,z,-,}
\end{tikzpicture}
\end{center}
\explique{Images des racines et des bornes :}{les nombres de la ligne des $x$}
\begin{calculs}
f(-1) &= 1+3-1 = 3 &&\justif{$-(-1)^{3}=+1$}\\
f(0) &= -1 &&\\
f(2) &= -8+12-1 = 3 &&\\
f(3) &= -27+27-1 = -1 &&
\end{calculs}
\begin{center}
\begin{tikzpicture}[scale=0.9]
\tkzTabInit[lgt=2,espcl=2.2]{$x$/0.8,$f'(x)$/0.8,$f(x)$/1.6}{$-1$,$0$,$2$,$3$}
\tkzTabLine{,-,z,+,z,-,}
\tkzTabVar{+/$3$,-/$-1$,+/$3$,-/$-1$}
\end{tikzpicture}
\end{center}
\cas{4.}
\explique{$f'$ s'annule en changeant de signe en $0$ ($-$ puis $+$) et en $2$ ($+$ puis
$-$) : $f$ a un \res{minimum local $-1$ en $0$} et un \res{maximum local $3$ en
$2$}.}{théorème 2}
\explique{Sur $\intff{-1}{3}$, le maximum est $\res{3}$ (atteint en $-1$ et en $2$) et le
minimum $\res{-1}$ (atteint en $0$ et en $3$).}{on compare les sommets, bornes
comprises}
\cas{5.}
\begin{calculs}
f(1) &= -1+3-1 = 1 &&\\
f'(1) &= -3+6 = 3 &&\justif{avec $f'(x)=-3x^{2}+6x$}\\
y &= f'(1)(x-1)+f(1) &&\justif{propriété 1}\\
y &= 3(x-1)+1 &&\\
y &= \res{3x-2} &&
\end{calculs}
\cas{6.}
\explique{Tableau de valeurs, au pas de $0{,}5$ :}{calculatrice}
\begin{center}
\renewcommand{\arraystretch}{1.3}
\begin{tabular}{|c|c|c|c|c|c|c|c|c|c|}
\hline
$x$ & $-1$ & $-0{,}5$ & $0$ & $0{,}5$ & $1$ & $1{,}5$ & $2$ & $2{,}5$ & $3$\\ \hline
$f(x)$ & $3$ & $-0{,}125$ & $-1$ & $-0{,}375$ & $1$ & $2{,}375$ & $3$ & $2{,}125$ & $-1$\\ \hline
\end{tabular}
\end{center}
\explique{On place les points du tableau, les extremums (tangentes horizontales en $0$ et en
$2$), puis $T$, qui passe par $(1\,;1)$ et $(0\,;-2)$.}{deux points suffisent pour tracer
une droite}
\begin{center}
\begin{tikzpicture}[x=1cm,y=1cm]
  \repere{-2}{-2}{4}{4}
  \gradx{-1,1,2,3} \grady{-1,1,2,3}
  \draw[coulCourbe,line width=1.1pt,domain=-1:3,samples=100]
    plot (\x,{-\x*\x*\x+3*\x*\x-1});
  \draw[coulExo,line width=0.9pt] (0,-2) -- (1.8,3.4);
  \draw[coulExo,line width=0.9pt] (-0.6,-1) -- (0.6,-1) (1.4,3) -- (2.6,3);
  \path[coulCourbe] (0,-1) node[croix]{} (2,3) node[croix]{} (1,1) node[croix]{};
  \node[coulExo,font=\footnotesize] at (1.95,3.75) {$T$};
  \node[coulCourbe,font=\footnotesize] at (3.4,-0.6) {$\Cf$};
\end{tikzpicture}
\end{center}
\rem{En $1$, la courbe traverse sa tangente : ce n'est pas une erreur de tracé.}
\end{exo}

\begin{exo}[Étude d'un quotient défini sur $\R$]
\cas{1.}
\explique{$x^{2}\ge 0$, donc $x^{2}+2\ge 2>0$ : le dénominateur ne s'annule jamais, $f$
est \res{définie sur $\R$}.}{on cherche les valeurs interdites}
\cas{2.}
\begin{calculs}
u(x) &= 2x+1 &&\justif{$u'(x)=2$}\\
v(x) &= x^{2}+2 &&\justif{$v'(x)=2x$}\\
f'(x) &= \dfrac{2(x^{2}+2)-(2x+1)\times 2x}{(x^{2}+2)^{2}} &&\justif{$\frac{u'v-uv'}{v^{2}}$}\\
f'(x) &= \dfrac{2x^{2}+4-4x^{2}-2x}{(x^{2}+2)^{2}} &&\\
f'(x) &= \dfrac{-2x^{2}-2x+4}{(x^{2}+2)^{2}} &&\\
f'(x) &= \dfrac{-2(x^{2}+x-2)}{(x^{2}+2)^{2}} &&\justif{on factorise par $-2$}\\
f'(x) &= \res{\dfrac{-2(x+2)(x-1)}{(x^{2}+2)^{2}}} &&\justif{$(x+2)(x-1)=x^{2}-x+2x-2$}
\end{calculs}
\cas{3.}
\explique{$(x^{2}+2)^{2}>0$ et $-2<0$ : on étudie le signe de chaque facteur.}{tableau de
signes}
\begin{center}
\begin{tikzpicture}[scale=0.9]
\tkzTabInit[lgt=2,espcl=2.2]{$x$/0.8,$-2$/0.8,$x+2$/0.8,$x-1$/0.8,$f'(x)$/0.8}{$-\infty$,$-2$,$1$,$+\infty$}
\tkzTabLine{,-,t,-,t,-,}
\tkzTabLine{,-,z,+,t,+,}
\tkzTabLine{,-,t,-,z,+,}
\tkzTabLine{,-,z,+,z,-,}
\end{tikzpicture}
\end{center}
\begin{calculs}
f(-2) &= \dfrac{-4+1}{4+2} = \dfrac{-3}{6} &&\\
f(-2) &= -\dfrac{1}{2} &&\\
f(1) &= \dfrac{2+1}{1+2} = 1 &&
\end{calculs}
\begin{center}
\begin{tikzpicture}[scale=0.9]
\tkzTabInit[lgt=2,espcl=2.2]{$x$/0.8,$f'(x)$/0.8,$f(x)$/1.6}{$-\infty$,$-2$,$1$,$+\infty$}
\tkzTabLine{,-,z,+,z,-,}
\tkzTabVar{+/{},-/$-\frac{1}{2}$,+/$1$,-/{}}
\end{tikzpicture}
\end{center}
\cas{4.}
\explique{$f'$ s'annule en changeant de signe : \res{minimum local $-\frac{1}{2}$ en
$-2$} et \res{maximum local $1$ en $1$}.}{théorème 2}
\cas{5.}
\begin{calculs}
f(0) &= \dfrac{1}{2} &&\\
f'(0) &= \dfrac{-2\times 2\times(-1)}{2^{2}} &&\justif{on remplace $x$ par $0$ dans la forme factorisée}\\
f'(0) &= \dfrac{4}{4} = 1 &&\\
y &= f'(0)(x-0)+f(0) &&\\
y &= \res{x+\dfrac{1}{2}} &&
\end{calculs}
\cas{6.}
\explique{Valeurs arrondies au centième :}{calculatrice}
\begin{center}
\renewcommand{\arraystretch}{1.3}
\begin{tabular}{|c|c|c|c|c|c|c|c|c|c|}
\hline
$x$ & $-4$ & $-3$ & $-2$ & $-1$ & $0$ & $1$ & $2$ & $3$ & $4$\\ \hline
$f(x)$ & $-0{,}39$ & $-0{,}45$ & $-0{,}5$ & $-0{,}33$ & $0{,}5$ & $1$ & $0{,}83$ & $0{,}64$ & $0{,}5$\\ \hline
\end{tabular}
\end{center}
\begin{center}
\begin{tikzpicture}[x=1.6cm,y=1.6cm]
  \repere{-4}{-1}{4}{2}
  \gradx{-4,-3,-2,-1,1,2,3,4} \grady{-1,1,2}
  \draw[coulCourbe,line width=1.1pt,domain=-4:4,samples=120]
    plot (\x,{(2*\x+1)/(\x*\x+2)});
  \draw[coulExo,line width=0.9pt] (-1.5,-1) -- (1.5,2);
  \draw[coulExo,line width=0.9pt] (-2.5,-0.5) -- (-1.5,-0.5) (0.5,1) -- (1.5,1);
  \path[coulCourbe] (-2,-0.5) node[croix]{} (1,1) node[croix]{} (0,0.5) node[croix]{};
  \node[coulExo,font=\footnotesize] at (1.7,1.85) {$T$};
  \node[coulCourbe,font=\footnotesize] at (3.8,0.75) {$\Cf$};
\end{tikzpicture}
\end{center}
\end{exo}

\begin{exo}[Étude d'un quotient avec une valeur interdite]
\cas{1.}
\explique{Le dénominateur s'annule pour $x-2=0$, soit $x=2$ : $f$ est définie sur
$\res{\R\setminus\{2\}}$.}{valeur interdite}
\cas{2.}
\begin{calculs}
u(x) &= x^{2}-3 &&\justif{$u'(x)=2x$}\\
v(x) &= x-2 &&\justif{$v'(x)=1$}\\
f'(x) &= \dfrac{2x(x-2)-(x^{2}-3)\times 1}{(x-2)^{2}} &&\justif{$\frac{u'v-uv'}{v^{2}}$}\\
f'(x) &= \dfrac{2x^{2}-4x-x^{2}+3}{(x-2)^{2}} &&\\
f'(x) &= \dfrac{x^{2}-4x+3}{(x-2)^{2}} &&\\
f'(x) &= \res{\dfrac{(x-1)(x-3)}{(x-2)^{2}}} &&\justif{$(x-1)(x-3)=x^{2}-3x-x+3$}
\end{calculs}
\cas{3.}
\explique{Pour $x\ne 2$, $(x-2)^{2}>0$ : $f'(x)$ a le signe de $(x-1)(x-3)$. La double
barre marque la valeur interdite $2$.}{tableau de signes}
\begin{center}
\begin{tikzpicture}[scale=0.9]
\tkzTabInit[lgt=2,espcl=2]{$x$/0.8,$x-1$/0.8,$x-3$/0.8,$f'(x)$/0.8}{$-\infty$,$1$,$2$,$3$,$+\infty$}
\tkzTabLine{,-,z,+,t,+,t,+,}
\tkzTabLine{,-,t,-,t,-,z,+,}
\tkzTabLine{,+,z,-,d,-,z,+,}
\end{tikzpicture}
\end{center}
\begin{calculs}
f(1) &= \dfrac{1-3}{1-2} = \dfrac{-2}{-1} &&\\
f(1) &= 2 &&\\
f(3) &= \dfrac{9-3}{3-2} = 6 &&
\end{calculs}
\begin{center}
\begin{tikzpicture}[scale=0.9]
\tkzTabInit[lgt=2,espcl=2]{$x$/0.8,$f'(x)$/0.8,$f(x)$/1.6}{$-\infty$,$1$,$2$,$3$,$+\infty$}
\tkzTabLine{,+,z,-,d,-,z,+,}
\tkzTabVar{-/{},+/$2$,-D+/{}/{},-/$6$,+/{}}
\end{tikzpicture}
\end{center}
\cas{4.}
\explique{\res{Maximum local $2$ en $1$} et \res{minimum local $6$ en $3$}.}{théorème 2}
\explique{Le maximum local est plus petit que le minimum local : c'est possible, car la
valeur interdite $2$ sépare la courbe en deux morceaux.}{chaque extremum ne vaut
qu'autour de son point}
\cas{5.}
\begin{calculs}
f(0) &= \dfrac{-3}{-2} = \dfrac{3}{2} &&\\
f'(0) &= \dfrac{(-1)\times(-3)}{(-2)^{2}} = \dfrac{3}{4} &&\\
y &= f'(0)(x-0)+f(0) &&\\
y &= \res{\dfrac{3}{4}x+\dfrac{3}{2}} &&
\end{calculs}
\cas{6.}
\explique{Valeurs arrondies au centième, de part et d'autre de la valeur interdite
:}{calculatrice}
\begin{center}
\renewcommand{\arraystretch}{1.3}
\begin{tabular}{|c|c|c|c|c|c|c|}
\hline
$x$ & $-3$ & $-2$ & $-1$ & $0$ & $1$ & $1{,}5$\\ \hline
$f(x)$ & $-1{,}2$ & $-0{,}25$ & $0{,}67$ & $1{,}5$ & $2$ & $1{,}5$\\ \hline
\end{tabular}
\quad
\begin{tabular}{|c|c|c|c|c|c|c|}
\hline
$x$ & $2{,}5$ & $3$ & $4$ & $5$ & $6$ & $7$\\ \hline
$f(x)$ & $6{,}5$ & $6$ & $6{,}5$ & $7{,}33$ & $8{,}25$ & $9{,}2$\\ \hline
\end{tabular}
\end{center}
\begin{center}
\begin{tikzpicture}[x=0.6cm,y=0.6cm]
  \repere{-4}{-2}{8}{10}
  \gradx{-3,-2,-1,1,2,3,4,5,6,7} \grady{-1,1,2,3,4,5,6,7,8,9}
  \draw[coulCourbe,line width=1.1pt,domain=-3:1.5,samples=80]
    plot (\x,{(\x*\x-3)/(\x-2)});
  \draw[coulCourbe,line width=1.1pt,domain=2.5:7,samples=80]
    plot (\x,{(\x*\x-3)/(\x-2)});
  \draw[coulExo,line width=0.9pt] (-3,-0.75) -- (2,3);
  \draw[coulExo,line width=0.9pt] (0.4,2) -- (1.6,2) (2.4,6) -- (3.6,6);
  \path[coulCourbe] (1,2) node[croix]{} (3,6) node[croix]{} (0,1.5) node[croix]{};
  \node[coulExo,font=\footnotesize] at (2.4,3.4) {$T$};
  \node[coulCourbe,font=\footnotesize] at (6.4,9.6) {$\Cf$};
\end{tikzpicture}
\end{center}
\rem{On ne relie jamais les deux morceaux de la courbe : $f$ n'existe pas en $2$.}
\end{exo}

\partie{Problèmes concrets}

\begin{exo}[Le vol d'une montgolfière]
\cas{1.}
\explique{La montgolfière décolle du village : l'altitude du village est celle de la
montgolfière à $t=0$.}{$f(0)$ : la valeur au départ}
\begin{calculs}
f(0) &= \dfrac{8\times 0+6}{0^{2}+1} = 6 &&
\end{calculs}
\explique{$f(t)$ est en centaines de mètres : le village est à $6\times 100=\res{600}$~m
d'altitude.}{on n'oublie pas l'unité}
\cas{2.}
\begin{calculs}
u(t) &= 8t+6 &&\justif{$u'(t)=8$}\\
v(t) &= t^{2}+1 &&\justif{$v'(t)=2t$}\\
f'(t) &= \dfrac{8(t^{2}+1)-(8t+6)\times 2t}{(t^{2}+1)^{2}} &&\justif{$\frac{u'v-uv'}{v^{2}}$}\\
f'(t) &= \dfrac{8t^{2}+8-16t^{2}-12t}{(t^{2}+1)^{2}} &&\\
f'(t) &= \dfrac{-8t^{2}-12t+8}{(t^{2}+1)^{2}} &&\\
f'(t) &= \dfrac{-4(2t^{2}+3t-2)}{(t^{2}+1)^{2}} &&\justif{on factorise par $-4$}\\
f'(t) &= \res{\dfrac{-4(2t-1)(t+2)}{(t^{2}+1)^{2}}} &&\justif{$(2t-1)(t+2)=2t^{2}+4t-t-2$}
\end{calculs}
\cas{3.}
\explique{Sur $\intff{0}{5}$, $(t^{2}+1)^{2}>0$ et $t+2>0$ ; comme $-4<0$, $f'(t)$ a le
signe contraire de $2t-1$.}{on ne garde que le facteur dont le signe change}
\explique{$2t-1$ s'annule en $t=0{,}5$, négatif avant, positif après : $f'(t)$ est positif
sur $\intfo{0}{0{,}5}$ et négatif sur $\intof{0{,}5}{5}$.}{signe contraire}
\begin{calculs}
f(0{,}5) &= \dfrac{4+6}{0{,}25+1} = \dfrac{10}{1{,}25} &&\\
f(0{,}5) &= 8 &&\\
f(5) &= \dfrac{46}{26} = \dfrac{23}{13} &&\justif{environ $1{,}8$}
\end{calculs}
\begin{center}
\begin{tikzpicture}[scale=0.9]
\tkzTabInit[lgt=2,espcl=2.4]{$t$/0.8,$f'(t)$/0.8,$f(t)$/1.6}{$0$,$0{,}5$,$5$}
\tkzTabLine{,+,z,-,}
\tkzTabVar{-/$6$,+/$8$,-/$\frac{23}{13}$}
\end{tikzpicture}
\end{center}
\cas{4.}
\explique{La montgolfière perd de l'altitude dès que $f$ décroît, donc à partir de
$t=0{,}5$.}{maximum de $f$}
\begin{calculs}
0{,}5~\text{h} &= 0{,}5\times 60~\text{min} &&\\
0{,}5~\text{h} &= 30~\text{min} &&
\end{calculs}
\explique{Elle commence à descendre à 10~h~50 $+$ 30~min, soit à $\res{11~\text{h}~20}$, à
l'altitude $f(0{,}5)=8$ centaines de mètres, soit $\res{800}$~m.}{on répond par une phrase,
avec l'unité}
\cas{5.}
\explique{On résout $f(t)=6$ (altitude du village), avec $t>0$ : $t=0$, le décollage, est
déjà une solution.}{on traduit la question}
\begin{calculs}
\dfrac{8t+6}{t^{2}+1} &= 6 &&\\
8t+6 &= 6(t^{2}+1) &&\justif{on multiplie par $t^{2}+1>0$}\\
8t+6 &= 6t^{2}+6 &&\\
6t^{2}-8t &= 0 &&\justif{on regroupe tout dans un membre}\\
2t(3t-4) &= 0 &&\justif{on factorise par $2t$}
\end{calculs}
\explique{Produit nul : $t=0$ (le décollage) ou $3t-4=0$, soit $t=\dfrac{4}{3}$.}{équation
produit nul}
\begin{calculs}
\dfrac{4}{3}~\text{h} &= 1~\text{h}+\dfrac{1}{3}\times 60~\text{min} &&\justif{$\frac{4}{3}=1+\frac{1}{3}$}\\
\dfrac{4}{3}~\text{h} &= 1~\text{h}~20~\text{min} &&
\end{calculs}
\explique{Elle revient à l'altitude du village à 10~h~50 $+$ 1~h~20, soit à
$\res{12~\text{h}~10}$.}{}
\end{exo}

\begin{exo}[Des bactéries et un antibiotique]
\cas{1.}
\begin{calculs}
N(0) &= 0-0+0+10 = 10 &&
\end{calculs}
\explique{$N(t)$ est en milliers : il y a $\res{10\,000}$ bactéries au départ.}{$N(0)$ : la
valeur au départ}
\cas{2.}
\begin{calculs}
N'(t) &= 3t^{2}-12\times 2t+36 &&\justif{on dérive ; $10$ donne $0$}\\
N'(t) &= 3t^{2}-24t+36 &&
\end{calculs}
\explique{Pour montrer une forme factorisée, le plus simple est de la développer et de
retrouver $N'(t)$.}{pas besoin de $\Delta$}
\begin{calculs}
3(t-2)(t-6) &= 3(t^{2}-6t-2t+12) &&\\
3(t-2)(t-6) &= 3(t^{2}-8t+12) &&\\
3(t-2)(t-6) &= 3t^{2}-24t+36 &&\justif{on retrouve $N'(t)$}
\end{calculs}
\cas{3.}
\explique{$N'(t)$ est un trinôme de racines $2$ et $6$, avec $a=3>0$ : positif à
l'extérieur des racines, négatif entre.}{signe du trinôme}
\begin{calculs}
N(2) &= 8-48+72+10 = 42 &&\\
N(6) &= 216-432+216+10 = 10 &&\\
N(8) &= 512-768+288+10 = 42 &&
\end{calculs}
\begin{center}
\begin{tikzpicture}[scale=0.9]
\tkzTabInit[lgt=2,espcl=2.2]{$t$/0.8,$N'(t)$/0.8,$N(t)$/1.6}{$0$,$2$,$6$,$8$}
\tkzTabLine{,+,z,-,z,+,}
\tkzTabVar{-/$10$,+/$42$,-/$10$,+/$42$}
\end{tikzpicture}
\end{center}
\cas{4.}
\begin{calculs}
N'(0) &= 3\times(-2)\times(-6) = 36 &&\justif{forme factorisée}\\
N'(4) &= 3\times 2\times(-2) = -12 &&
\end{calculs}
\explique{$N'(0)=36>0$ : au départ, la population augmente à la vitesse de $\res{36\,000}$ bactéries par
heure.}{$N'(t)$ : vitesse, en milliers de bactéries par heure}
\explique{$N'(4)=-12<0$ : à $t=4$, la population diminue à la vitesse de $\res{12\,000}$ bactéries par
heure.}{une vitesse négative : la grandeur diminue}
\cas{5.}
\begin{calculs}
N(t)-10 &= t^{3}-12t^{2}+36t &&\justif{on retire $10$}\\
t(t-6)^{2} &= t(t^{2}-12t+36) &&\justif{$(a-b)^{2}=a^{2}-2ab+b^{2}$}\\
t(t-6)^{2} &= t^{3}-12t^{2}+36t &&\justif{on retrouve $N(t)-10$}
\end{calculs}
\explique{Sur $\intff{0}{8}$, $t\ge 0$ et $(t-6)^{2}\ge 0$ : $N(t)-10\ge 0$, donc
$N(t)\ge 10$.}{un produit de deux nombres positifs est positif}
\explique{\res{Non} : le nombre de bactéries ne descend jamais sous $10\,000$. Il y revient
exactement à $t=6$, puis remonte.}{$N(6)=10$ : le minimum du tableau}
\end{exo}

\begin{exo}[Le coût moyen d'un artisan]
\cas{1.}
\begin{calculs}
M(10) &= 10+\dfrac{400}{10} = 50 &&
\end{calculs}
\explique{Pour $10$ lampes fabriquées, une lampe coûte en moyenne
$\res{50~\text{\euro}}$.}{on interprète, avec l'unité}
\cas{2.}
\begin{calculs}
M'(x) &= 1+400\times\left(-\dfrac{1}{x^{2}}\right) &&\justif{$\frac{400}{x}=400\times\frac{1}{x}$}\\
M'(x) &= \dfrac{1\fois{x^{2}}}{1\fois{x^{2}}}-\dfrac{400}{x^{2}} &&\justif{même dénominateur}\\
M'(x) &= \dfrac{x^{2}-400}{x^{2}} &&\\
M'(x) &= \res{\dfrac{(x-20)(x+20)}{x^{2}}} &&\justif{$a^{2}-b^{2}=(a-b)(a+b)$, avec $b=20$}
\end{calculs}
\cas{3.}
\explique{Sur $\intff{5}{50}$, $x^{2}>0$ et $x+20>0$ : $M'(x)$ a le signe de $x-20$,
négatif avant $20$, positif après.}{on ne garde que le facteur dont le signe change}
\begin{calculs}
M(5) &= 5+80 = 85 &&\\
M(20) &= 20+20 = 40 &&\\
M(50) &= 50+8 = 58 &&
\end{calculs}
\begin{center}
\begin{tikzpicture}[scale=0.9]
\tkzTabInit[lgt=2,espcl=2.4]{$x$/0.8,$M'(x)$/0.8,$M(x)$/1.6}{$5$,$20$,$50$}
\tkzTabLine{,-,z,+,}
\tkzTabVar{+/$85$,-/$40$,+/$58$}
\end{tikzpicture}
\end{center}
\explique{Le coût moyen est le plus bas pour $\res{20}$ lampes ; il vaut alors
$\res{40~\text{\euro}}$.}{minimum de $M$ sur $\intff{5}{50}$}
\cas{4.}
\explique{On utilise les variations : $M(10)=50$ (question 1) et $M(40)=40+10=50$.}{on
cherche où $M$ revient à $50$}
\explique{Sur $\intff{10}{20}$, $M$ décroît de $50$ à $40$ ; sur $\intff{20}{40}$, elle
croît de $40$ à $50$. Donc $M(x)\le 50$ pour tout $x$ de $\intff{10}{40}$ : l'affirmation
est \res{vraie}.}{une fonction monotone reste entre ses valeurs aux bornes}
\rem{Autre méthode : résoudre $M(x)\le 50$ ; en multipliant par $x>0$, on obtient
$x^{2}-50x+400\le 0$, trinôme de racines $10$ et $40$, négatif entre elles.}
\cas{5.}
\explique{Deux droites sont parallèles quand elles ont le même coefficient directeur. Celui
de la tangente au point d'abscisse $10$ est $M'(10)$ ; celui de la droite est
$-2$.}{coefficient directeur de la tangente : $M'(a)$}
\begin{calculs}
M'(10) &= \dfrac{(10-20)(10+20)}{10^{2}} &&\\
M'(10) &= \dfrac{-10\times 30}{100} = -3 &&
\end{calculs}
\explique{$-3\ne -2$ : la tangente n'est pas parallèle à la droite d'équation $y=-2x+1$.
L'affirmation est \res{fausse}.}{}
\end{exo}

\end{document}
